% Modernised reading — fonds Alexandre Grothendieck, shelfmark 136, pages 1-375.
%
% This is an INTERPRETATION, not a transcription. It re-reads the pages in
% today's notation and vocabulary, states the hypotheses the manuscript leaves
% implicit, and drops the critical apparatus so that the mathematics reads
% straight through. Everything the brackets and underlines carried — a doubtful
% reading, a notation he used differently, a missing passage — moves into a
% footnote.
%
% It is held to being mathematically correct as it stands, not merely faithful
% to the page. Where the manuscript is loose or mistaken, the true statement is
% what appears here, and the footnote says what the page has. In any dispute of
% substance the transcription governs: whoever cites Grothendieck must cite the
% batch-NN.fr.tex files.
%
% Scope: ONE document for the whole shelfmark, nineteen batches read as one
% dossier. The folder is not one argument but a bundle of strands around the
% 1976 IHES lecture; the spine section after the résumé names them and their
% page ranges. Pages 356-375 are notes on Steenrod operations, Eilenberg-MacLane
% spaces and Ext(G_a, G_a) whose author the pages do not name; they are read as
% notes, and attributed to nobody.
%
% NOTATION fixed for the whole folder:
%   sSet (his Ss, \widehat\Delta) the topos of simplicial sets; X_* an object,
%     X_{[n]} its n-simplices. His « ½simplicial » is read « simplicial »
%     throughout (the old sense), footnoted once;
%   k a commutative ring, S_k = k{T} the divided-power polynomial algebra,
%     T^{(i)} its divided powers, c_{i,j} = (i+j)!/(i! j!);
%   \Phi_* = \Phi_{*k} : [n] -> k^{n+1} (= i_{0*}k); the pages also call it
%     \mathfrak{G}_* (pp. 101-107) and \mathfrak{L}_* (pp. 224-232);
%     k_* the constant object, \Psi_* = \Phi_*/k_*;
%   C^{p,q}_* = \Gamma^p\Phi_* (x) \Lambda^q\Psi_*, with p the COMPLEMENTARY
%     (divided-power) degree and q the EXTERIOR degree, total degree n = p+q.
%     The pages index by (p,q) on pp. 7-9 and by (total, exterior) on pp. 80
%     and 230-232; every statement here is converted to (p,q);
%   \mathcal{D}^{p,q}_* = \Gamma^p\Phi_* (x) \Lambda^q\Phi_*, the larger
%     resolution of k_* of pp. 16 and 73 (his D and C^{**});
%   \omega_N, \tau_N, \varphi_N the shadow, truncation and coreflection
%     functors on graded S-modules; \mathcal{M}_N the modules in degrees >= N;
%     \tau_\omega the passage to the quotient by negligible modules;
%   \mathcal{A}_*(S, S^+, t) the simplicial divided-power differential algebra
%     of pp. 5, 124, 161, 203; \Omega the group scheme of automorphisms (p. 177);
%   V(M) : A -> Hom_k(M, A) the « vector bundle » functor of a k-module.
%   M[i] is the shift (M[i])_j = M_{i+j}; the pages also write T for it
%     (p. 243) and t (p. 286), letters this reading keeps for other things.
%
% Substantive departures from the pages, each footnoted where it occurs:
%   - p. 44: the « largest divisible submodule » is not \bigcap nN in general;
%     the argument gives the largest p-divisible submodule, k a Z_(p)-algebra;
%   - p. 40: the last vanishing condition is c_{d,\beta-d-1} u = 0;
%   - p. 140: Dirichlet's integral makes the constant \gamma_{i} equal to 1
%     (ours), so the Stokes integral of a divided-power monomial over the
%     simplex of side t is t^{(n + \sum i)};
%   - p. 165: the Leibniz rule of the \delta_i carries binomial coefficients;
%   - p. 191: \varphi(dX_0) = dX_0 + df, not 1 + df;
%   - p. 214: the divided powers of the free alternating algebra live on the
%     even part (Cartan's convention); on odd elements they are consistent
%     only for strictly commuting pairs, as p. 216 itself finds;
%   - p. 221: the operator d\delta + \delta d kills cohomology degree by
%     TOTAL degree;
%   - p. 228: \Gamma^i\Psi_* and Sym^i\Psi_* are NOT K(k,i) in general
%     (the décalage of p. 361 gives LSym^i(k[1]) = 0 for i >= 2);
%   - p. 240: v_p(p^{(i)}) <= v_p(p^{(i+1)}) iff p^2 does not divide i+1;
%   - pp. 262-264: the vanishing of R^n\varphi_N on M (x) S is drawn from the
%     exactness of \tau i alone, which does not suffice; recorded as a gap;
%   - p. 364: Adem gives Sq^1 Sq^2 = Sq^3;
%   - p. 369: K_2(R) = H_2(E(R)) = \pi_2(BGL(R)^+), not \pi_1(E(R)); the
%     Bloch formula is in codimension i;
%   - p. 371: Cartan's symmetry holds in the stable range, with a degree shift.
%
% Announced in the folder and never established: that the divided-power De Rham
% complex with all its structures determines the homotopy type (pp. 11-12,
% 16-17); that the groupoid of admissible theories is connected and simply
% connected (p. 120); the second of the « two groups » of p. 105; the
% reconstruction of K^\bullet in D(C) rests on the gap of pp. 262-264; the
% splitting s of p. 77 is never defined; the Künneth formula (p. 13) is not
% developed; the Steenrod squares hoped for in characteristic 2 (p. 181) are
% not found.
%
% Third-party material that the transcriptions only note — a personal letter
% (p. 117), letters (pp. 119, 244), typescripts by other hands (pp. 215, 217,
% 348-352) — is not restated beyond saying that it is there (RIGHTS.md).
%
% Pass: Opus 5.5 (claude-opus-5-5), 2026-10-10 — first pass, unchecked against
% the pages by a human.
\documentclass[11pt,a4paper]{article}
\input{../preamble/grothendieck}

\folder{136}
\pages{1}{375}
\dating{[à partir de 1975-1976]}
\watermark{Édition de démonstration\\interprétation personnelle de l'œuvre}
\foldertitle{Complexe de De Rham à puissance divisée [conférence de 1976 à l’IHÉS] : notes manuscrites (s.d.) — lecture modernisée du dossier entier}

\begin{document}

\section*{Résumé}

\begin{resume}
Trois cent soixante-quinze feuillets, réunis sous une couverture orange qui
porte « Complexe DR à p.d.\ et ombres », préparent et prolongent une conférence
donnée à l'IHÉS en 1976. La question est simple à dire : peut-on écrire, pour
un espace quelconque, des « formes différentielles » qui calculent sa
cohomologie à coefficients \emph{entiers} — sans perdre la torsion ?

Depuis De Rham, les formes différentielles calculent la cohomologie, mais à
coefficients réels ; Sullivan, au début des années 1970, avait montré qu'en
remplaçant les variétés par des simplexes et les fonctions lisses par des
polynômes à coefficients rationnels, on récupère même tout le type d'homotopie
rationnel d'un espace. Mais tout cela vit en caractéristique nulle, là où les
nombres $2, 3, 5, \ldots$ sont inversibles et où la torsion — le fait qu'un
cycle puisse être nul une fois répété $p$ fois sans l'être lui-même — devient
invisible. « Du point de vue de Weil, c'était un défaut impardonnable »,
note-t-il dans la marge de la première page. L'obstacle est presque scolaire :
intégrer $x^n\,dx$ demande de diviser par $n+1$.

L'idée vient de la géométrie algébrique en caractéristique $p$, où elle avait
déjà servi : les \emph{puissances divisées}. Plutôt que d'écrire $x^n$ et de
diviser ensuite, on prend pour symbole de base $x^{(n)}$, « $x^n/n!$ » — comme
on tient un registre des coefficients d'une série de Taylor plutôt que des
dérivées. Les primitives redeviennent entières. Mais un nouvel obstacle
surgit : le simplexe est défini par l'équation $x_0 + \cdots + x_n = 1$, et le
nombre $1$ n'admet pas de puissances divisées entières. Sa solution est de
remplacer $1$ par un paramètre $T$ qui, lui, en admet : on travaille avec des
simplexes « de taille $T$ ». Le prix est qu'on ne récupère plus la cohomologie
elle-même, mais ce qu'il appelle ses \emph{ombres} : des copies de la
cohomologie multipliée par les puissances divisées de $T$, dont on ne voit que
la partie située au-dessus d'une certaine hauteur. Tout un pan du dossier
établit, à coups d'arithmétique des coefficients binomiaux, qu'une ombre suffit
à reconstituer l'objet qui la projette : vue d'assez haut, la silhouette
détermine la forme.

Autour de ce noyau, le dossier construit les outils et pose les questions.
L'intégration de Stokes compare ces formes aux cochaînes singulières, et
l'intégrale d'un monôme en puissances divisées sur le simplexe de taille $t$
se révèle être une puissance divisée de $t$. Il cherche si la structure est
unique, et calcule pour cela le groupe des automorphismes de l'algèbre
simpliciale universelle : un groupe multiplicatif prolongé par un groupe
unipotent, qui ne produit aucune opération cohomologique nouvelle dès que $2$
est inversible — d'où l'espoir, non réalisé, de voir apparaître en
caractéristique $2$ les carrés de Steenrod. Il dualise tout en termes de
coalgèbres, met en place les catégories de modules gradués où vivent les
ombres, et termine sur une réflexion de méthode.

Cette réflexion (pages 301 à 317) vaut d'être lue pour elle-même. Les invariants
qu'on attache naïvement à un espace sont, dit-il, « trop gros et frustes » : ils
dépendent de la façon dont l'espace est découpé, pas seulement de sa forme. Le
travail consiste à les munir de structures supplémentaires raisonnables, puis à
« inverser les flèches » qui viennent des équivalences d'homotopie — et
l'invariant obtenu, plus grossier, est aussi plus subtil, et « moins
redondant ». Il ajoute un aveu d'atelier : des deux langages possibles, objets
simpliciaux et complexes de chaînes, le second est plus économique, mais seul le
premier laisse voir les produits et les puissances extérieures, qui
« s'introduisent de façon impérieuse ». Les vingt dernières pages, d'une autre
facture, sont des notes sur les opérations de Steenrod, les espaces
d'Eilenberg–MacLane et les extensions du groupe additif.

Ce qu'on cherchera aujourd'hui sous d'autres noms : les formes polynomiales de
Sullivan et l'homotopie rationnelle ; le complexe de De Rham à puissances
divisées de la cohomologie cristalline ; la correspondance de Dold–Kan ; le
théorème de Mandell, qui récupère le type d'homotopie à partir des cochaînes
munies de leur structure $E_\infty$ — la réponse moderne, par un autre chemin,
à la question que ce dossier pose.

\keywords{divided power de Rham complex, polynomial de Rham forms, Sullivan model, rational homotopy theory, integral cochain models, Mandell theorem, divided powers, Lucas theorem, Kummer theorem, Legendre formula, graded modules, truncation functor, Serre quotient category, localization of abelian categories, simplicial sets, presheaf topos, Stokes integration map, normalized cochains, Dirichlet integral, Koszul complex, Poincare lemma, Cartier isomorphism, flasque resolution, simplicial commutative ring, automorphism group scheme, Hasse-Schmidt derivation, coalgebra, comodule, rational module, cotensor product, affine scheme of algebras, free divided power algebra, Dold-Kan correspondence, Dold-Puppe derived functors, decalage, Eilenberg-MacLane object, Kan extension, nerve and realization, free cocompletion, abelianization, Eilenberg-Zilber theorem, Alexander-Whitney map, shuffle map, Steenrod operations, Steenrod algebra, Adem relations, Milnor dual Steenrod algebra, cohomology of Eilenberg-MacLane spaces, Dold-Thom theorem, bar construction, Milnor K2, Steinberg group, Gersten resolution, Bloch formula, Ext of the additive group, perfect site, uniformization, projective structure, Euler class}
\end{resume}

\section*{Le fil du dossier, et l'ordre des feuillets}

Le dossier n'est pas une rédaction suivie. C'est une liasse de notes — séries
paginées par l'auteur (pages 2 à 17, 301 à 317, 361 à 365), feuilles isolées,
chemises de sa main (« Gribouillis vrac », « Au fil des jours… », page 318 ;
« Gribouillis », page 335) — dont une bonne part est écrite au verso de
tapuscrits sans rapport (épreuves de SGA~7, textes sur les accouplements de
monodromie et les modèles de Néron, listings), que les transcriptions laissent
de côté. L'ordre archivistique n'est pas l'ordre de rédaction : la version
correcte du théorème sur les homomorphismes de l'algèbre simpliciale (page 161)
précède dans la liasse la version fausse qu'il corrige (page 193).

\subsection*{Les stations}

\begin{enumerate}
\item \emph{Le plan de la conférence} (pages 2 à 17) : historique ; la
construction $C^{\bullet\bullet}_{\mathrm{DRpd}}(X_*,k)$ ; le calcul de sa
cohomologie ; les ombres ; la liste des problèmes ouverts.
\item \emph{L'arithmétique des ombres} (pages 19 à 44) : critère de Lucas,
pleine fidélité des foncteurs « ombres » $\omega_N$, le théorème sur
$M \otimes_k S$, les homomorphismes homogènes entre tronqués.
\item \emph{Cochaînes et intégration de Stokes} (pages 47 à 77) : cochaînes des
ensembles simpliciaux, intégrales de Stokes, comparaison avec De Rham
classique, le complexe vu comme noyau de la multiplication par $dT$.
\item \emph{Sur un topos quelconque ; unicité} (pages 79 à 120) : faisceaux
tests, « condition de De Rham », algèbres de De Rham spéciales, théories
admissibles.
\item \emph{La version duale} (pages 122 à 159) : coalgèbres, schémas affines
en algèbres, chaînes de De Rham, dualisation de Stokes, le yoga coalgèbres
$\mapsto$ algèbres, comodules.
\item \emph{Homomorphismes et automorphismes} (pages 161 à 211) : le théorème
des pages 161 à 171, le schéma en groupes $\Omega$, l'espoir de la
caractéristique $2$, les endomorphismes de $\mathcal{A}_1$.
\item \emph{Outils} (pages 214 à 240) : algèbre alternée libre à puissances
divisées, les quatre complexes (De Rham, Koszul, et leurs versions à
puissances divisées), calculs d'homotopie, version complétée et filtrée.
\item \emph{Les ombres dans les catégories dérivées} (pages 243 à 298) :
catégories $\mathcal{M}_N$ et leurs adjoints, modules négligeables,
reconstruction d'un complexe à partir de son ombre, produit cotensoriel.
\item \emph{Algèbre universelle des invariants} (pages 300 à 354) :
conventions simpliciales, foncteurs universels, Dold–Puppe ; brouillons sur
coalgèbres et puissances divisées ; un feuillet étranger sur les courbes
(pages 325 à 327).
\item \emph{Notes sur les opérations cohomologiques} (pages 356 à 375) :
Eilenberg–Zilber, Steenrod, Eilenberg–MacLane, $K_2$ de Milnor,
$\operatorname{Ext}(\mathbf{G}_a, \mathbf{G}_a)$.
\end{enumerate}

\subsection*{Conventions, valables pour tout le dossier}

$k$ est un anneau commutatif, $S = S_k = k\{T\}$ l'algèbre des polynômes à
puissances divisées en une variable, de base les $T^{(i)}$ avec
$T^{(i)} T^{(j)} = c_{i,j}\, T^{(i+j)}$, où
$c_{i,j} = \frac{(i+j)!}{i!\,j!}$. Les ensembles simpliciaux forment le topos
$\mathrm{sSet} = \widehat{\Delta}$\footnote{Il écrit « ½simplicial » et
« ½cosimplicial », au sens ancien où « semi-simplicial » désignait ce qu'on
appelle aujourd'hui \emph{simplicial} (avec dégénérescences) ; c'est le sens
qu'il précise lui-même pages 3, 6 et 301. Nous écrivons « simplicial ».} ; on
note $X_*$ un objet et $X_{[n]}$ ses $n$-simplexes. On note
$\Phi_* : [n] \mapsto k^{n+1}$ le $k$-module simplicial des fonctions sur les
sommets, $k_* \subset \Phi_*$ le sous-module diagonal, $\Psi_* = \Phi_*/k_*$, et
$T \in \Gamma(\Phi_*)$ la section « constante égale à $1$ »\footnote{Le même
objet porte trois noms dans la liasse : $\Phi_*$ (pages 7, 73, 99, 336),
$\mathfrak{G}_*$ (pages 101 à 107, lecture incertaine de la lettre gothique) et
$\mathfrak{L}_*$ (pages 224 à 232). C'est l'image directe $i_{0*}(k)$ par le
point $i_0$ du topos (page 99).}.

Le complexe central est la $k\{T\}$-algèbre bigraduée simpliciale
$C^{p,q}_* = \Gamma^p_k \Phi_* \otimes_k \Lambda^q_k \Psi_*$, où $p$ est le
\emph{degré complémentaire} (le degré en puissances divisées), $q$ le
\emph{degré extérieur}, $n = p + q$ le degré total\footnote{Les pages changent
d'indexation : $(p,q)$ = (complémentaire, extérieur) pages 7 à 9, (total,
extérieur) pages 80 et 230 à 232, et (total, extérieur) encore pour les
structures de De Rham des pages 93 et suivantes. Tout est ramené ici à
(complémentaire, extérieur), et dit en clair quand une page indexe autrement.}.
Le complexe plus gros $\mathcal{D}^{p,q}_* = \Gamma^p \Phi_* \otimes \Lambda^q
\Phi_*$ des pages 16 et 73 est une résolution de $k_*$ qui le contient. La
translation des modules gradués est notée $M \mapsto M[i]$,
$(M[i])_j = M_{i+j}$\footnote{Il la note $T$ page 243 et $t$ page 286 ; ces
deux lettres désignent ailleurs la variable de $k\{T\}$ et l'élément marqué
d'un anneau à puissances divisées.}.

Trois homonymies à garder en tête. La lettre $K$ désigne le bicomplexe
$K^{\bullet\bullet}_u$ d'une immersion (pages 79 à 86), un foncteur des pages
341 à 344, et les objets d'Eilenberg–MacLane $K(G,n)$ ; ce sont trois objets
sans rapport. $\Phi_*$ (le module) n'a rien à voir avec $\varphi_N$ (le
foncteur adjoint de la troncature). Et l'« algèbre de De Rham standard »
$\mathrm{DR}(\Phi, T)$ des pages 95 à 107 est le même objet que
$C^{\bullet\bullet}_*$, vu sur un topos quelconque.

\subsection*{Ce qui n'est pas établi}

La question qui motive la conférence — le complexe de De Rham à puissances
divisées, avec toutes ses structures, détermine-t-il le type d'homotopie d'un
espace simplement connexe de type fini ? (pages 11 et 12) — reste posée ;
des modèles minimaux sur $\mathbb{Z}$ qui permettraient de l'aborder, le
dossier dit lui-même n'avoir « rien fait dans cette direction » (page 13).
Sont également annoncés sans être établis : la connexité et la simple connexité
du groupoïde des théories admissibles (page 120) ; le second des « deux
groupes » de la page 105 ; le scindage canonique cherché page 77 ; la formule
de Künneth (page 13). La reconstruction d'un complexe à partir de son ombre
(pages 262 à 272) repose sur une annulation d'$\operatorname{Ext}$ dont la
justification écrite ne suffit pas.

\pagerange{2}{3}
\section*{Le plan de la conférence : historique (pages 2 et 3)}

La série paginée de 1 à 16 par l'auteur s'ouvre sur un plan, encadré au crayon
« Plan conf.\ IHÉS en 1976 », avec en marge la liste des cadres où la question
se pose : géométrie différentielle, analytique, algébrique, arithmétique,
topologie algébrique (linéaire par morceaux, simpliciale).

\begin{enumerate}
\item[a)] Formes différentielles (Poincaré) et formule de Stokes
$\int_C d\omega = \int_{\partial C} \omega$.
\item[b)] Le théorème de De Rham\footnote{La page dit « conjecturé par
E.~Cartan » ; c'est son attribution, que nous reproduisons sans la
discuter.}, $H^*_{\mathrm{DR}}(X) \simeq H^*(X, \mathbb{C})$ ; on y perd la
torsion. Il est « maintenant bien compris » par la théorie des faisceaux et le
lemme de Poincaré.
\item[c)] La théorie de Hodge : bigraduation de $H^*_{\mathrm{DR}}(X)$ pour $X$
kählérienne compacte.
\item[d)] Le théorème de H.~Cartan et Serre sur les variétés de Stein, où le
lemme de Poincaré vaut dans le cadre analytique.
\item[e)] Les schémas sur un corps de caractéristique $p$ : Dwork,
Washnitzer–Monsky, puis le « yoga cristallin » de Berthelot, Illusie, Messing,
Mazur\footnote{En marge, au crayon : Hartshorne, Herrera, Bloch (lecture
incertaine), Ogus.}. On y trouve des cohomologies dont les coefficients ne
contiennent plus $\mathbb{Q}$, et où la torsion n'est en principe pas perdue —
« du point de vue de Weil », note la marge, la perte de la torsion « était un
défaut impardonnable ».
\item[f)] Son propre théorème de comparaison pour les variétés algébriques sur
$\mathbb{C}$, généralisé à des coefficients par Deligne et Hartshorne, qui donne
en caractéristique nulle la signification topologique de la cohomologie de De
Rham algébrique.
\item[g)] Le complexe de De Rham–Sullivan d'un espace topologique : formes
différentielles polynomiales sur les simplexes singuliers, à coefficients dans
$\mathbb{Q}$, $\mathbb{R}$ ou $\mathbb{C}$, qui calcule la cohomologie à ces
coefficients. Sullivan montre davantage : le foncteur $C^\bullet_{\mathrm{DRS}}$
induit une équivalence entre la catégorie homotopique rationnelle des espaces
connexes et simplement connexes de type fini et une catégorie dérivée de
$\mathbb{Q}$-algèbres différentielles graduées commutatives (au sens gradué) de
degrés $\geq 0$, avec $H^0 \simeq \mathbb{Q}$ et $H^1 = 0$\footnote{L'hypothèse
de type fini (sur les groupes d'homotopie rationnels, ou la cohomologie) est
ajoutée par nous ; elle est nécessaire à l'énoncé de Sullivan, et la page dit
« sauf erreur ». Un mot illisible suit « simplement connexes », qui la portait
peut-être. Le nom moderne est la \emph{théorie de l'homotopie rationnelle} ; le
foncteur est celui des \emph{formes polynomiales} $A_{\mathrm{PL}}$, et
l'équivalence est due à Quillen et Sullivan.}. La théorie des modèles minimaux
lit les $\pi_i \otimes \mathbb{Q}$ sur un tel modèle. Mais on perd à nouveau la
torsion.
\end{enumerate}

\pagerange{3}{5}
\section*{La démonstration de Sullivan, et l'obstacle sur $\mathbb{Z}$ (pages 3 à 5)}

Soit $E^{[n]}$ l'espace affine $\sum_0^n x_i = 1$ de $\mathbb{R}^{n+1}$, avec sa
$\mathbb{Q}$-structure, et
$\mathrm{DRS}^\bullet_{[n]} = C^\bullet_{\mathbb{Q}\text{-DR}}(E^{[n]})$ son
complexe de De Rham algébrique sur $\mathbb{Q}$ :
\[
\mathrm{DRS}^\bullet_{[n]} = \mathbb{Q}[X_i, dX_i]_{0 \leq i \leq n} \big/ \Bigl(\textstyle\sum X_i - 1,\ \sum dX_i\Bigr) \simeq \mathbb{Q}[X_0, \ldots, X_{n-1}][dX_0, \ldots, dX_{n-1}] .
\]
C'est contravariant en $[n]$ : on obtient une $\mathbb{Q}$-algèbre
différentielle graduée commutative simpliciale $\mathrm{DRS}^\bullet_*$. Pour
un ensemble simplicial $X_*$ (par exemple le complexe singulier $S_*(X)$ d'un
espace),
\[
C^\bullet_{\mathrm{DRS}}(X_*, \mathbb{Q}) = \operatorname{Hom}_{\mathrm{sSet}}(X_*, \mathrm{DRS}^\bullet_*) ,
\]
et le foncteur ainsi défini sur $\mathrm{sSet}$ est représenté par
$\mathrm{DRS}^\bullet_*$. Le théorème d'isomorphisme
$H^q(C^\bullet_{\mathrm{DRS}}(X_*)) \simeq H^q(X_*, \mathbb{Q})$ résulte de deux
faits : $\mathrm{DRS}^\bullet_*$ est une résolution de $\mathbb{Q}_*$ (lemme de
Poincaré algébrique sur les espaces affines rationnels) ; et chaque composante
$\mathrm{DRS}^i_*$ est un objet acyclique du topos $\mathrm{sSet}$, ce qui
revient à dire que ses groupes d'homotopie sont nuls. Le même argument vaut
pour des formes $C^\infty$, analytiques réelles ou complexes.

Sur $\mathbb{Z}$, c'est le premier fait qui tombe : pour intégrer $x^n\,dx$ il
faut le dénominateur de $x^{n+1}/(n+1)$\footnote{Un passage de deux lignes et
demie, biffé et en grande partie illisible, disait qu'on « aimerait trouver une
$\mathbb{Z}$-algèbre différentielle » analogue, « mais il paraît qu'il y a des
[\ldots] contre ». Le fait aujourd'hui connu est qu'il n'existe pas de modèle
fonctoriel strictement commutatif des cochaînes à coefficients entiers, ni même
dans $\mathbb{F}_p$ : les opérations de Steenrod en sont l'obstruction. Nous ne
savons pas ce que portait le mot illisible.}. La géométrie algébrique connaît
le moyen de « sauter à pieds joints par-dessus le canular » : les puissances
divisées. Si $E^{[n]}$ avait une origine sur $\mathbb{Z}$, on prendrait le
complexe de De Rham à puissances divisées d'un $\mathbb{Z}$-module libre. Mais
il n'en a pas : dans $\mathbb{Z}\{X_0, \ldots, X_n\}$, l'élément
$\sum X_i - 1$ n'appartient pas à l'idéal à puissances divisées, et l'on ne peut
« diviser par l'idéal à p.d.\ engendré » par lui.

\pagerange{5}{8}
\section*{La construction (pages 5 à 8)}

\subsection*{Le simplexe de taille $t$}

On remplace l'anneau de coefficients $\mathbb{Z}$ par un anneau $S$ muni d'un
idéal à puissances divisées $J$ et d'un « paramètre » $t \in J$, et l'équation
$\sum X_i = 1$ par $\sum X_i = t$. On pose
\[
C^\bullet_{\mathrm{DRpd}[n]}(S, J, t) = S\{X_0, \ldots, X_n\}[dX_0, \ldots, dX_n] \big/ \Bigl(\textstyle\sum X_i - t,\ \sum dX_i\Bigr)_{\mathrm{pd}} ,
\]
où l'on divise par l'idéal à puissances divisées engendré par
$\sum X_i - t$ et $\sum dX_i$ ; cela a un sens, car l'idéal des formes
d'augmentation dans $J$ est un idéal à puissances divisées. C'est une
$S$-algèbre différentielle graduée alternée, de degrés $\geq 0$, augmentée vers
$S/J$, à puissances divisées sur l'idéal d'augmentation et compatibles à la
différentielle ($d(x^{(n)}) = x^{(n-1)}\,dx$), et elle est isomorphe à
$S\{X_0, \ldots, X_{n-1}\}[dX_0, \ldots, dX_{n-1}]$ : c'est une
\emph{résolution de $S$}. Pour $[n]$ variable on obtient une algèbre
simpliciale $C^\bullet_{\mathrm{DRpd}*}(S,J,t)$, fonctorielle en $(S,J,t)$, et
pour un ensemble simplicial $X_*$
\[
C^\bullet_{\mathrm{DRpd}}(X_*; S, J, t) = \operatorname{Hom}\bigl(X_*, C^\bullet_{\mathrm{DRpd}*}(S,J,t)\bigr) ,
\]
contravariant en $X_*$. On a seulement un homomorphisme
$H^\bullet_{\mathrm{DRpd}}(X_*; S,J,t) \to H^\bullet(X_*, S)$, qui en général
n'est pas un isomorphisme.

\subsection*{Le cas universel $S = k\{T\}$}

On prend $S = k\{T\}$, $J = k\{T\}^+$, $t = T$. Alors
$k\{T, X_0, \ldots, X_n\}/(\sum X_i - T)_{\mathrm{pd}} \simeq k\{X_0, \ldots,
X_n\}$, et
\[
C^{\bullet\bullet}_{\mathrm{DRpd}*}(k) \overset{\mathrm{def}}{=} C^\bullet_{\mathrm{DRpd}*}\bigl(k\{T\}, k\{T\}^+, T\bigr) \simeq \Gamma^\bullet_k \Phi_* \otimes_k \Lambda^\bullet \Psi_* .
\]
C'est une algèbre \emph{bigraduée} — degré complémentaire $p$, degré
extérieur $q$ —, unitaire, associative, alternée pour le degré total
(anticommutative, et de carré nul en degré impair), augmentée vers $k$, à
puissances divisées sur l'idéal d'augmentation, avec une différentielle de
bidegré $(-1, +1)$, donc de degré total nul, compatible aux puissances
divisées. Ces structures passent aux
\[
C^{\bullet\bullet}_{\mathrm{DRpd}}(X_*, k) = \operatorname{Hom}\bigl(X_*, C^{\bullet\bullet}_{\mathrm{DRpd}*}(k)\bigr) ,
\]
contravariant en $X_*$, covariant en $k$.

\subsection*{La cohomologie}

$\Phi_*$ est homotope à $0$ (c'est la proposition 1 de la page 226) ; donc
$C^{p,q}_* = \Gamma^p \Phi_* \otimes \Lambda^q \Psi_*$ est homotope à $0$, et
acyclique, pour $p \neq 0$. En degré total $n$ fixé, la ligne
\[
0 \to k_* \xrightarrow{\ 1 \mapsto T^{(n)}\ } C^{n,0}_* \to C^{n-1,1}_* \to \cdots \to C^{1,n-1}_* \to C^{0,n}_* \to 0
\]
est exacte\footnote{La page parle de « résolution de longueur $n$ » sans
justifier l'exactitude ; c'est l'exactitude du complexe de Koszul à puissances
divisées associé à l'extension $0 \to k_* \to \Phi_* \to \Psi_* \to 0$, $\Psi_*$
étant plat (cf.\ les pages 79 et 218). L'extrémité droite de la suite est
d'une lecture incertaine sur la page.}, et tous ses termes sont acycliques sauf
le dernier : c'est le tronqué en degré $n$ d'une résolution acyclique de $k_*$.
D'où, pour tout ensemble simplicial $X_*$,
\[
H^{p,q}_{\mathrm{DRpd}}(X_*, k) \simeq H^q(X_*, k) \qquad \text{pour tout } p \geq 0 ,
\]
et $0$ pour $p < 0$\footnote{La page écrit la condition « $q \leq p+q$,
i.e.\ $p \geq 0$ », lecture incertaine ; c'est bien ce que donne la troncature.
La même conclusion est retrouvée page 232, en indexation (total, extérieur) :
$H^{n,q} \simeq H^q$ pour $q \leq n$.}.

\pagerange{9}{11}
\section*{Les ombres (pages 9 à 11)}

Pour $q$ fixé, la structure de $k\{T\}$-module de $H^{\bullet, q}$ est la
suivante : la classe d'un $h \in H^q(X_*, k)$ dans la ligne de degré total $n$
correspond à l'augmentation $1 \mapsto T^{(n)}$, et la multiplication par
$T^{(j)}$ envoie la ligne $n$ sur la ligne $n+j$ par $c_{j,n}$. Indexé par le
degré total, on trouve donc
\[
H^{\bullet, q}_{\mathrm{DRpd}}(X_*, k) \simeq \tau_{\geq q}\bigl(H^q(X_*, k) \otimes_k k\{T\}\bigr) ,
\]
le module gradué $H^q \otimes k\{T\}$ tronqué en degrés $\geq q$ : c'est ce
qu'il appelle une \emph{$q$-ombre}. Les isomorphismes sont compatibles aux
structures multiplicatives. A priori on ne récupère donc pas les $H^q$ et leurs
cup-produits, mais seulement leurs ombres et les accouplements qu'elles
héritent.

Ce qui donne l'espoir que le complexe à puissances divisées porte une
information homotopique précise, c'est l'observation suivante, démontrée aux
pages 21 à 34 : un $k$-module $M$ se reconstitue, à isomorphisme canonique près,
à partir de \emph{n'importe lequel} de ses tronqués $\tau_{\geq q}(M \otimes_k
S)$, et un accouplement $M \otimes N \to P$ à partir de l'accouplement
correspondant des tronqués, pour $q$ assez grand. Cette « théorie des ombres »
acquise, la connaissance de $C^{\bullet\bullet}_{\mathrm{DRpd}}(X_*, k)$
comme $S$-algèbre différentielle bigraduée implique celle de l'algèbre graduée
$H^\bullet(X_*, k)$ — et, en raffinant un peu, celle de
$\mathrm{R}\Gamma(X_*, k)$ dans $D(k)$ avec l'accouplement
$\mathrm{R}\Gamma \otimes^{\mathrm{L}} \mathrm{R}\Gamma \to \mathrm{R}\Gamma$
venant du cup-produit\footnote{Aucun nom moderne ne s'impose pour ces
« ombres » ; ce sont des modules gradués sur l'algèbre des puissances divisées
d'une variable, tronqués, et l'énoncé de reconstruction est celui des pages 21
à 34. La page 10 est d'une écriture très chargée, et la phrase sur
$\mathrm{R}\Gamma$ n'y est lue qu'en partie ; c'est l'objet des pages 262 à
272.}.

\pagerange{11}{17}
\section*{Remarques et problèmes (pages 11 à 17)}

\emph{a) Le lien avec Sullivan.} Si $k$ est une $\mathbb{Q}$-algèbre,
$C^\bullet_{\mathrm{DRS}}(X_*, k) \simeq C^{\bullet\bullet}_{\mathrm{DRpd}}(X_*,
k)/(T-1)$. Plus généralement, pour tout $k$, le quotient par $T - 1$ donne
$C^\bullet_{\mathrm{DRS}}(X_*, k \otimes \mathbb{Q})$ — ce qui s'explique par
$\mathbb{Z}\{T\}/(T-1) \simeq \mathbb{Q}$, $T^{(n)} \mapsto 1/n!$\footnote{Cette
explication est la nôtre. Le passage au quotient ne commute pas aux produits
infinis : l'énoncé est vrai tel quel pour $X_*$ n'ayant qu'un nombre fini de
simplexes non dégénérés, et nous ne l'avons pas vérifié au-delà. Pour un triple
$(S,J,t)$ général la page propose
$C^\bullet_{\mathrm{DRpd}}(X_*; S,J,t)/(t-1) \simeq C^\bullet_{\mathrm{DRS}}(X_*,
S_{\mathbb{Q}}/(t-1))$, avec un « ? » en marge et un membre de droite
surchargé.}.

\emph{b) La question principale.} Il a jusqu'ici négligé la structure à
puissances divisées sur l'idéal d'augmentation de
$C^{\bullet\bullet}_{\mathrm{DRpd}}(X_*, k) \to H^0(X_*, k)$, qui est
importante. Appelons \emph{complexe de De Rham à puissances divisées virtuel}
sur $k$ une $k$-algèbre bigraduée différentielle, augmentée, associative,
unitaire, alternée, à différentielle de bidegré $(-1,+1)$ et bidegrés
$\geq 0$, avec puissances divisées sur l'idéal d'augmentation compatibles à la
différentielle, et telle que les $H^{\bullet q}$ soient des $q$-ombres ;
inversons les quasi-isomorphismes. On obtient un foncteur de la catégorie
homotopique $(\mathrm{Hot})$ vers cette catégorie dérivée $(\mathrm{DRpd})$, et
l'on peut se demander si sa restriction aux espaces connexes et simplement
connexes à $H^i(X, \mathbb{Z})$ de type fini induit une équivalence avec les
complexes sur $\mathbb{Z}$ tels que $H^0 \simeq \mathbb{Z}$, $H^1 = 0$ et les
$H^i$ de type fini\footnote{La marge ajoute « avec conditions de finitude sur
$\mathcal{X}_*$ et sur les $\pi_i(\mathcal{X}_*)$ ». La question est restée
ouverte dans le dossier. La réponse moderne passe par un autre chemin : Mandell
(2006) a montré que les cochaînes singulières entières, munies de leur
structure d'algèbre $E_\infty$, déterminent le type d'homotopie des espaces
nilpotents de type fini. Rien n'indique que le manuscrit ait connu cette
direction, et nous ne prétendons pas que la structure à puissances divisées en
soit l'équivalent.}.

\emph{c) Opérations cohomologiques.} Il a cherché à reconstruire les
opérations de type Steenrod ou Whitney à partir de
$C^{\bullet\bullet}_{\mathrm{DRpd}}(X_*, -)$, et n'a que des résultats partiels
négatifs : avec les seules opérations multiplicatives de
$C^{\bullet\bullet}_{\mathrm{DRpd}*}(k)$, on ne trouve rien d'intéressant. (Les
pages 173 à 181 précisent ce résultat négatif.)

\emph{d) Modèles minimaux.} Il faudrait étudier des modèles minimaux à la
Sullivan sur $\mathbb{Z}$. Il n'a « rien fait dans cette direction », ni même
développé une formule de Künneth pour un produit $X_* \times Y_*$ : la
difficulté technique est qu'on ne peut sans doute pas supposer les complexes
plats sur $k\{T\}$.

\emph{e) La version homologique.} De même que le complexe de chaînes
$C_\bullet(X_*, \mathbb{Z})$ est l'objet « le plus fin », dont on déduit toutes
les cochaînes par $C^\bullet(X_*, k) \simeq \operatorname{Hom}^\bullet_{\mathbb{Z}}
(C_\bullet(X_*, \mathbb{Z}), k)$, il existe une cogèbre différentielle
bigraduée $C^{\mathrm{DRpd}}_{\bullet\bullet}(X_*, k)$, à comultiplication
divisée et coaugmentée, telle que $C^{\bullet\bullet}_{\mathrm{DRpd}}(X_*, k')
\simeq \operatorname{Hom}_k(C^{\mathrm{DRpd}}_{\bullet\bullet}(X_*, k), k')$
pour toute $k$-algèbre $k'$, et
$C^{\mathrm{DRpd}}_{\bullet\bullet}(X_*, k') \simeq
C^{\mathrm{DRpd}}_{\bullet\bullet}(X_*, k) \otimes_k k'$. C'est
$C^{\mathrm{DRpd}}_{\bullet\bullet}(X_*, \mathbb{Z})$ qu'il faudrait considérer
pour aborder b), c), d) sans conditions de finitude ; il commute aux limites
inductives quelconques en $X_*$. (Construction aux pages 122 à 134.)

\emph{g) Faisceautisation.} La définition a un sens pour un Anneau
commutatif $k$ d'un topos. On aimerait définir des $\mathrm{R}f_*$ dans des
catégories dérivées convenables, pour $f$ un morphisme de topos, par exemple en
remplaçant un complexe par un complexe quasi-isomorphe à composantes flasques —
mais y en a-t-il toujours ? Le même problème se pose déjà, sur $\mathbb{Q}$,
pour les complexes de De Rham–Sullivan : c'est un problème posé par Deligne
pour $\widehat{\Delta} \to$ (point) ; « sauf erreur », il y a un topos qui
marche pour les espaces topologiques\footnote{La lettre f) est biffée et
remplacée par g), puis g) par h) au paragraphe suivant ; il n'y a pas de f)
sur la page.}.

\emph{h) Des invariants linéaires plus fins.} Le complexe
$C^{\bullet\bullet}_{\mathrm{DRpd}*}(k)$ se déduit de
$\mathcal{D}^{\bullet\bullet}_* = \Gamma^\bullet \Phi_* \otimes
\Lambda^\bullet \Phi_*$ — une algèbre différentielle à puissances divisées qui
est une résolution de $k$ — et de $dT$ : en chaque bidegré,
\[
C^{p,q}_* \simeq \operatorname{Ker}\bigl(\mathcal{D}^{p,q+1}_* \xrightarrow{\ dT \cdot\ } \mathcal{D}^{p,q+2}_*\bigr) ,
\]
d'où la même formule pour les sections sur $X_*$, et les structures de
$C^{\bullet\bullet}_{\mathrm{DRpd}}$ peuvent se lire comme déduites de
structures (compliquées) de $\mathcal{D}^{\bullet\bullet}(X_*, k)$. On peut
aussi définir « \emph{la} structure multiplicative à puissances divisées
cohomologique du type d'homotopie $X_*$ », qui permettrait de reconstituer
aussi bien $C_{\mathrm{DRpd}}$ que $\mathcal{D}$ ou les complexes de
Sullivan. Il ignore si le complexe à puissances divisées suffit à la récupérer ;
sinon c'est cette dernière qui serait la candidate pour exprimer le type
d'homotopie. (La construction est reprise page 73.)

\pagerange{19}{22}
\section*{Coefficients binomiaux modulo $p$ (pages 19 à 22)}

\textbf{Lemme 1 (Lucas).} \emph{Soient $p$ premier, $i = \sum_\varepsilon
u_\varepsilon p^\varepsilon$ et $j = \sum_\varepsilon v_\varepsilon
p^\varepsilon$ les développements $p$-adiques, $0 \leq u_\varepsilon,
v_\varepsilon \leq p-1$. Alors $c_{i,j} \not\equiv 0 \pmod p$ si et seulement si
$u_\varepsilon + v_\varepsilon \leq p-1$ pour tout $\varepsilon$, et en tout
cas $c_{i,j} \equiv \prod_\varepsilon c_{u_\varepsilon, v_\varepsilon} \pmod
p$.}\footnote{Les bornes des sommes sont écrites $e$, celles du produit
$\nu$ ; l'anneau du calcul est écrit $\mathbb{F}_p[V,Y]$ pour
$\mathbb{F}_p[X,Y]$ ; et la conclusion de la première étape est écrite
« $c_{v,s} \not\equiv 0$ » pour $c_{r,s}$. Ce sont les théorèmes de Lucas (1878)
et, pour le critère, de Kummer (1852) ; la page ne les nomme pas.}

On écrit $i = u + pr$, $j = v + ps$, on développe $(X+Y)^{i+j} =
(X^p+Y^p)^{r+s}(X+Y)^{u+v}$ dans $\mathbb{F}_p[X,Y]$, et l'on obtient
$c_{i,j} \equiv c_{u,v}\, c_{r,s}$, avec $c_{u,v} \equiv 0$ exactement quand
$u + v \geq p$ ; le lemme suit par récurrence sur la longueur des
développements.

\textbf{Corollaire 2.} \emph{Si $N + 1 = (N'+1)p$ et $\varepsilon \geq 1$, alors
$c_{N+1, p^\varepsilon - 1} \equiv c_{N'+1, p^{\varepsilon-1}-1}$, qui vaut $1$
si $p^{\varepsilon-1} \mid N'+1$ et $0$ sinon ; et $c_{N, p^\varepsilon} \equiv
c_{N', p^{\varepsilon-1}} \equiv v_{\varepsilon-1} + 1$, où $N' = \sum
v_\alpha p^\alpha$.} C'est le décalage d'un chiffre entre les développements
de $N$ et de $N'$.

\textbf{Lemme 2.} \emph{Soient $E$ un $\mathbb{Z}_{(p)}$-module, $N \in
\mathbb{N}$, $e$ tel que $p^e \geq N+1$, et des $x_i \in E$ pour $i \in \{N\}
\cup \{N + p^\varepsilon\}_{0 \leq \varepsilon \leq e}$, tels que
$\binom{j}{i}(x_j - x_i) = 0$ pour $i \in \{N, N+1\}$ et $i < j$. Alors $x_N =
x_{N+1}$.} On pose $y_\varepsilon = x_{N+p^\varepsilon} - x_N$ ; on déduit la
conclusion $y_0 = 0$ du système plus faible obtenu en ne gardant que les
équations dont le coefficient est une unité, par récurrence sur $N$ : si $p
\nmid N+1$ c'est immédiat, et sinon le corollaire 2 ramène le système pour $N$
au même système pour $N' = \frac{N+1}{p} - 1 < N$ (« cqfd. Ouf ! »). La marge
remarque que sur un $\mathbb{F}_p$-module on ne peut pas conclure que tous les
$x_i$ sont égaux.

\textbf{Théorème 1.} \emph{Soient $E$ un groupe abélien, $N \in \mathbb{N}$, et
$(x_i)_{i \in \mathbb{N}}$ dans $E$ tels que $\binom{j}{i}(x_j - x_i) = 0$ pour
$j > i \geq N$. Alors $x_i = x_N$ pour tout $i \geq N$.} On localise en chaque
premier et l'on applique le lemme 2 de proche en proche.

\pagerange{22}{25}
\section*{Les foncteurs « ombres » $\omega_N$ (pages 22 à 25)}

Soient $S^{+(N)} = \sum_{i \geq N} k\, T^{(i)}$ et, pour un $k$-module $M$,
le $S$-module gradué $\omega_N(M) = M \otimes_k S^{+(N)}$\footnote{La page écrit
$M \otimes_S S^{+(N)}$ ; le produit tensoriel est sur $k$.}.

\textbf{Corollaire 1.} \emph{a) Le foncteur $\omega_N$, des $k$-modules vers les
$S$-modules gradués, est pleinement fidèle. b) L'application $x \mapsto (x\,
T^{(i)})_{i \geq N}$ identifie $M$ aux familles $(x_i)_{i \geq N}$, $x_i \in M\,
T^{(i)}$, telles que $T^{(j)} x_i = c_{i,j}\, x_{i+j}$ pour $i \geq N$, $j \geq
1$.} Écrivant $x_i = \xi_i T^{(i)}$, la condition devient $c_{i,j}(\xi_{i+j} -
\xi_i) = 0$, et le théorème 1 conclut ; pour a), un morphisme
$\omega_N(M) \to \omega_N(M')$ est une famille $u_i : M \to M'$ avec
$c_{i,j}(u_{i+j} - u_i) = 0$, d'où $u_i$ constant.

\emph{Remarque 1.} $\omega_N$ a un adjoint à droite $\pi_N$, avec
$\pi_N(\mathcal{M}) = \{(x_i)_{i \geq N} \mid x_i T^{(j)} = c_{i,j} x_{i+j}\}
\subset \prod_{i \geq N} \mathcal{M}_i$, et b) dit que l'unité $M \to
\pi_N \omega_N(M)$ est un isomorphisme.

\emph{Remarque 2.} $\omega_N(M) = S^{+(N)} \cdot \omega_0(M)$ est la troncature
$\tau_{\geq N}$ de $\omega_{N'}(M)$ pour $N \geq N'$. Les ombres sont donc
d'autant plus grossières que $N$ est grand, et toutes pleinement fidèles —
alors que la troncature $\tau_N$ elle-même ne l'est pas.

Soit $\mathrm{Omb}(k)$ la catégorie des « ombres virtuelles », les
$S$-modules gradués. Les modules nuls en degrés $\geq N$ forment une
sous-catégorie épaisse, et le quotient (les « $N$-ombres virtuelles ») est
canoniquement équivalent aux modules nuls en degrés $< N$, le foncteur quotient
s'identifiant à $\tau_N$\footnote{Le signe entre $i$ et $N$ dans la définition
est surchargé ; la lecture $\mathcal{M}_i = 0$ pour $i \geq N$ est celle
qu'impose la suite.}. Les modules \emph{négligeables} — ceux dont chaque élément
est annulé par $S^{+(N)}$ pour $N$ assez grand — forment encore une
sous-catégorie épaisse ; la phrase sur le localisé s'interrompt, et la notion
est reprise pages 253 à 258.

\pagerange{26}{34}
\section*{Le théorème sur $M \otimes_k S$ (pages 26 à 34)}

\textbf{Théorème.} \emph{Pour tout $k$-module $M$ et tout $N \geq 0$,
l'homomorphisme canonique $M \otimes_k S \to \varphi_N \tau_N(M \otimes_k S)$
est un isomorphisme en degrés $\geq 0$ ; de façon équivalente, pour $0 \leq i
\leq N$,}
\[
(M \otimes_k S)_i = M\, T^{(i)} \xrightarrow{\ \sim\ } \operatorname{Hom}^i\bigl(\tau_{N-i}(S), M \otimes_k S\bigr) .
\]
Ici $\varphi_N$ est l'adjoint à droite de la troncature (pages 247 et 290) ; en
degrés $\geq N$ c'est clair, et le cas $i = 0$ est le corollaire 1. Un
homomorphisme de degré $i$ de $\tau_{N-i}(S)$ dans $M \otimes S$ s'écrit
$T^{(k-i)} \mapsto \xi_k T^{(k)}$ pour $k \geq N$, et la linéarité s'écrit
$c_{l,k}\,\xi_k = c_{l,k-i}\,\xi_{k+l}$. Le théorème équivaut donc à l'énoncé
où $k$ ne figure plus :

\textbf{Théorème.} \emph{Soient $0 \leq i \leq N$, $M$ un groupe abélien et
$(\xi_k)_{k \geq N}$ dans $M$ tels que $c_{l,k}\,\xi_k = c_{l,k-i}\,\xi_{k+l}$
pour $k \geq N$, $l \geq 0$. Il existe un unique $\xi \in M$ tel que $\xi_k =
\binom{k}{i}\,\xi$ pour tout $k \geq N$.}\footnote{Il écrit tantôt
$c_{k,k-i}$, tantôt $c_{i,k-i}$ ; ces deux coefficients sont égaux à
$\binom{k}{i}$.}

\emph{Démonstration.} On se ramène à un $\mathbb{Z}_{(p)}$-module : les
$\binom{k}{i}$, $k \geq N$, sont premiers entre eux dans leur ensemble (pour
chaque $p$, $k = i + p^\varepsilon$ avec $p^\varepsilon > i$ donne une unité),
d'où une combinaison $\sum m_\alpha \binom{k_\alpha}{i} = 1$ qui fixe $\xi$ et
son unicité, et l'égalité se vérifie dans chaque localisé. Sur
$\mathbb{Z}_{(p)}$, on retranche la solution évidente et l'on est ramené au
\textbf{Corollaire} : \emph{si $(\xi_k)$ satisfait la relation et s'annule en un
$k_0$ avec $\binom{k_0}{i}$ inversible, elle est nulle.} Cela se fait en deux
temps.

\textbf{Lemme 1.1.} \emph{Si $k' \geq k \geq N$, $j = k - i$, $j' = k'-i$, et
$c_{i,j}$, $c_{i,j'}$ sont des unités, alors $c_{j'-j, j}$ est une unité si et
seulement si $c_{j'-j, j+i}$ l'est, et dans ce cas $\xi_k = 0 \iff \xi_{k'} =
0$.} Les quotients $c_{j'-j,j+i}/c_{j'-j,j}$ et $c_{i,j'}/c_{i,j}$ sont égaux
dans $\mathbb{Q}$, et la relation avec $l = j'-j$ conclut.

\textbf{Lemme 1.} \emph{Les $\xi_k$ tels que $c_{i,k-i}$ soit une unité sont
tous nuls ou tous non nuls.} On complète $i + j$ jusqu'à $p^{r+1} - 1$ en
ajoutant chiffre à chiffre $\delta_\alpha = (p-1) - (i_\alpha + j_\alpha)$ ; la
marge en donne une présentation plus simple : pour $p^r > i+j$ et $\delta =
p^r - 1 - (i+j)$, on a $c_{i, j+\delta}$ et $c_{i+j, \delta}$ non nuls modulo
$p$, d'où $\xi_k = 0 \iff \xi_{p^r - 1} = 0$.

\textbf{Lemme 2.} \emph{Pour $i, j \geq 0$, il existe $l \geq 0$ tel que
$c_{i+j,l}$ et $c_{i,j+l}$ soient non nuls modulo $p$} : on prend $l = p^r - 1 -
(i+j)$ avec $p^r > i+j$\footnote{Une première démonstration par les retenues
de l'addition $p$-adique (pages 32 et 33) est barrée de longs traits obliques ;
la démonstration retenue est celle-ci.}. Le lemme commun, dit-il, est que
$\mu + \nu = p^r - 1$ entraîne $c_{\mu,\nu} \not\equiv 0 \pmod p$ — aucune
retenue n'est possible.

Le théorème annonce encore un corollaire, $\operatorname{Ext}^n\bigl(
\tau_{N-i}(S)[-i], M \otimes_k S\bigr) = 0$ pour $n > 0$, dont la lecture est
incertaine et dont la démonstration ne vient qu'aux pages 262 à 264 (voir
là-bas la réserve qu'elle appelle).

Le feuillet isolé de la page 34 consigne la formule de Legendre
$v_p(n!) = \frac{n - s_p(n)}{p-1}$, où $s_p(n)$ est la somme des chiffres de $n$
en base $p$, d'où $v_p(p^n/n!) = \frac{s_p(n) + n(p-2)}{p-1}$, et le calcul
$v_p\bigl(\frac{(p^{r+1})!}{(p^r!)^p\, p!}\bigr) = 0$, qui sert page 42.

\pagerange{36}{44}
\section*{Homomorphismes homogènes entre tronqués (pages 36 à 44)}

Pour $N^*$ dans la catégorie $\mathcal{M}_N$ des modules gradués nuls en
degrés $< N$, l'adjonction donne
\[
\operatorname{Hom}^d(M^*, \varphi_N N^*) \simeq \operatorname{Hom}^d(\tau_{N-d} M^*, N^*) ,
\]
et en particulier, pour des $k$-modules $M$, $N$ et $d \geq 0$,
$\operatorname{Hom}^d(M \otimes_k S, N \otimes_k S) \simeq
\operatorname{Hom}^d(\tau_{N-d} M_S, \tau_N N_S) \simeq \operatorname{Hom}_k(M,
N)$\footnote{Il emploie la même lettre $N$ pour l'entier de troncature et pour
le $k$-module. Le bas de la page 36, écrit tête-bêche et barré, porte un calcul
sans rapport (« Diagramme de Cayley »).}. On étudie ensuite les homomorphismes
de degré $d$ de $\tau_\alpha M_S$ dans $\tau_\beta N_S$, $\beta \geq \alpha$,
où $M_S = M \otimes_k S$.

a) Si $d \geq \beta - \alpha$, ils s'identifient à $\operatorname{Hom}_k(M, N)$,
$u$ correspondant à $m\,T^{(j)} \mapsto c_{d,j}\, u(m)\, T^{(d+j)}$.

b) Si $0 \leq d < \beta - \alpha$, ce sont les $u \in \operatorname{Hom}_k(M,
N)$ tels que $c_{d,j}\, u = 0$ pour $\alpha \leq j \leq \beta - d - 1$ : il faut
que les termes tombant en degré $d + j < \beta$ s'annulent\footnote{La page
écrit le système jusqu'à $c_{d,\beta-d}\,u = 0$, puis la condition « $\alpha
\leq j < \beta - d$ » ; la dernière équation doit être $c_{d, \beta-d-1}\,u =
0$, ce que dit la condition.}.

c) Si $d < 0$ — homomorphismes de degré négatif — on se place, au feuillet
de la page 42, sur une $\mathbb{Z}_{(p)}$-algèbre $k$. Posant $T_r =
T^{(p^r)}$, on a $T_r^{\,p} = p\,\nu'_r\, T_{r+1}$ avec $\nu'_r = (p-1)!\,
\frac{(p^{r+1})!}{(p^r!)^p\, p!}$ inversible dans $\mathbb{Z}_{(p)}$ (page 34) ;
$\tau_{p^r}S$ est engendré par les $T_s$, $s \geq r$, avec les relations
$T_s^{p-1} T_s = p\nu'_s T_{s+1}$. Un homomorphisme de degré $-d$, $0 < d <
p^r$, est donné par des $\eta_s = u(T_s) \in M \otimes S_{p^s - d}$, et en
écrivant $\eta_s = \xi_s\, T_\nu^{\delta_\nu} \cdots
T_{r-1}^{\delta_{r-1}}(T_r \cdots T_{s-1})^{p-1}$ ($\nu = v_p(d)$) la relation
devient $\xi_s = p\,\nu'_s\,\xi_{s+1}$. Chaque $\xi_s$ est donc le premier terme
d'une suite de divisions successives par $p$.

\textbf{Proposition.} \emph{Soient $k$ une $\mathbb{Z}_{(p)}$-algèbre, $M$, $N$
des $k$-modules, $d > 0$, $\alpha \geq d$, et $N_0$ le plus grand sous-module
$p$-divisible de $N$. Tout $u \in \operatorname{Hom}^{-d}(\tau_\alpha M_S, N_S)$
se factorise par $N_{0\,S}$ : $\operatorname{Hom}^{-d}(\tau_\alpha M_S,
N_{0\,S}) \xrightarrow{\sim} \operatorname{Hom}^{-d}(\tau_\alpha M_S,
N_S)$.}\footnote{Trois écarts. (i) La page définit $N_0 = \bigcap_{n} nN$ et
l'appelle « le plus grand sous-module divisible » ; ces deux objets diffèrent
en général ($\bigcap nN$ peut être strictement plus gros, ce qu'on appelle le
premier sous-groupe d'Ulm). Ce que l'argument donne, puisque chaque $\xi_s$ se
prolonge en une suite $\xi_s, \xi_{s+1}, \ldots$ de divisions successives, est
l'appartenance au plus grand sous-module $p$-divisible, et c'est ce que nous
énonçons. (ii) L'énoncé de la page vaut pour $k$ quelconque ; le calcul n'est
fait que pour $k$ une $\mathbb{Z}_{(p)}$-algèbre, comme le dit la marge de la
page 42. (iii) Les deux membres de l'isomorphisme sont écrits identiques sur la
page.} \emph{Corollaire : si $N$ n'a pas de sous-module divisible non nul, par
exemple si $k$ est de torsion, ces homomorphismes sont nuls.}

\pagerange{47}{51}
\section*{Cohomologie des ensembles simpliciaux (pages 47 à 51)}

Ces trois pages sont au crayon très pâle, et la transcription n'en lit guère
que les formules : la lecture est mince ici. On y voit le cadre suivant. Un
foncteur $C^*$ des ensembles simpliciaux vers les complexes de $k$-modules,
calculant fonctoriellement $H^*(X_\bullet, M)$ et dont chaque composante
transforme limites inductives en limites projectives, est représentable :
$C^n(X_\bullet) = \operatorname{Hom}(X_\bullet, C^n_\bullet)$, et les
$C^*_\bullet$ forment une résolution de $M$, dont les cycles $Z^n$ sont liés aux
objets d'Eilenberg–MacLane $K(M,n)$. Exemple 1 : les cochaînes ordinaires,
représentées par $C^n_{\bullet M}$ avec $(C^n_{\bullet M})_m = M^{\operatorname{Hom}
(\Delta_n, \Delta_m)}$. Exemple 2 : les cochaînes normalisées, représentées par
\[
(C^{n!}_{\bullet M})_m \simeq M^{\operatorname{Hom}_{\mathrm{mono}}(\Delta_n, \Delta_m)} \simeq \Lambda^{n+1}(k^{\Delta_m}) \otimes_k M ,
\]
les applications strictement croissantes $[n] \to [m]$ correspondant aux
parties à $n+1$ éléments. Cette identification des cochaînes normalisées aux
puissances extérieures de $\Phi_*$ revient pages 67 et 337.

\pagerange{53}{55}
\section*{Intégration des algèbres simpliciales (pages 53 à 55)}

Soit $A_*$ un groupe différentiel gradué simplicial ($d$ de degré $+1$), avec
$A^p_n = 0$ pour $p \notin [0,n]$, et tel que, muni de l'augmentation
$\varepsilon_n : A_0 \to A^*_n$, chaque $A^*_n$ soit une résolution de $A_0$.
Soit $\partial = \sum_i (-1)^i \partial^*_{n,i} : A^*_n \to A^*_{n-1}$
l'opérateur de bord des faces.

\textbf{Proposition 1.} \emph{Il existe un et un seul système d'homomorphismes
$K_n : A^*_n \to A_0$ tel que a) $K_0 = \mathrm{id}$ ; b) $K_n(A^p_n) = 0$ pour
$p \neq n$ ; c) (Stokes) $K_n(d\omega) = K_{n-1}(\partial\omega)$.}

Par récurrence : tout $\varpi \in A^N_N$ est un cobord $d\omega$ puisque
$H^N(A^*_N) = 0$, et l'on pose $K_N(d\omega) = K_{N-1}(\partial\omega)$ ; il
faut que cela ne dépende pas de $\omega$, c'est-à-dire que $d\omega = 0$
entraîne $K_{N-1}(\partial\omega) = 0$. Pour $N \geq 2$, $\omega = d\alpha$ et
$K_{N-1}(\partial d\alpha) = K_{N-1}(d\partial\alpha) = K_{N-2}(\partial\partial
\alpha) = 0$ ; pour $N = 1$, $\omega = \varepsilon_1(\alpha)$ et
$\partial\varepsilon_1(\alpha) = 0$.

\emph{Remarques.} 1) On n'a utilisé que $H^N(A^*_N) = 0$ pour $N \geq 1$,
$H^{N-1}(A^*_N) = 0$ pour $N \geq 2$, et $H^0(A^*_1) = \varepsilon_1(A_0)$.
2) Il y a une version abstraite, pour un bicomplexe augmenté satisfaisant ces
conditions de degré dans une catégorie abélienne. 3) Par unicité, ces
« intégrales » $A^n_n \to A_0$ commutent aux structures supplémentaires — la
fin de la remarque est illisible.

\pagerange{57}{59}
\section*{Le topos induit par un ensemble simplicial (pages 57 et 59)}

Pour $X \in \mathrm{sSet}$, le topos induit est $\mathrm{sSet}/X \simeq
\widehat{\Delta/X}$ ; l'étude cohomologique de $X$ est celle de ce topos. Un
faisceau abélien sur $X$ est un foncteur $(\Delta/X)^\circ \to \mathrm{Ab}$ ;
un objet simplicial $A_\bullet$ définit le faisceau $A \times X \to X$, composé
$(\Delta/X)^\circ \to \Delta^\circ \xrightarrow{A} \mathrm{Ens}$, et l'on peut
parler de l'hypercohomologie $\mathrm{R}\Gamma(X, A^*_\bullet)$ d'un objet
gradué. Le complexe des cochaînes d'un faisceau abélien $A$ est
$C^n(X, A) = \prod_{\sigma \in X_{[n]}} A(\sigma)$.

\pagerange{61}{68}
\section*{Cochaînes et systèmes de Stokes (pages 61 à 68)}

Un système de coefficients contravariant sur le simplexe variable
$A_* : (\Delta/X)^\circ \to \mathcal{E}$ donne les chaînes
$C_*(X, A_*) = i_!(A_*)$, où $i : \Delta/X \to \Delta$ ; un « cofaisceau »
covariant donne les cochaînes ; on passe de l'un à l'autre en remplaçant
$\mathcal{E}$ par $\mathcal{E}^\circ$\footnote{La moitié supérieure de la page
61, biffée, esquisse les mêmes formules pour une fibration $C \to B$.}.

Pour un $k$-Module $F$ sur $X_*$, on a $C^*(X_*, F) \simeq \Gamma(X_*,
C^*_{X_*} \otimes_{k_*} F)$, et $C^*_{X_*} \otimes F$ est une résolution
flasque de $F$ (« à vérifier en passant », dit la marge), que $F$ soit
localement constant ou non. D'où $H^*(C^*(X_*, F)) \simeq H^*(X_*, F)$, et
de même pour l'hypercohomologie d'un complexe, $\mathrm{R}\Gamma(X_*, F^*)
\simeq C^*(X_*, F^*)$ (complexe simple associé).

Soit maintenant $A^*_*$ un complexe de $k$-modules simpliciaux, et $M$ un
$k$-module. Un morphisme de complexes $A^*_* \to C^*_{*M}$ (les cochaînes
représentantes) est un système $k_n : A^n_n \to M$ satisfaisant la
\emph{condition de Stokes} $k_n \circ \partial_{n+1} = k_{n+1} \circ d$ sur
$A^n_{n+1}$. Il induit, pour tout $X_*$, un morphisme $\Gamma(X_*, A^*_{X_*})
\to C^*(X_*, M)$, $f \mapsto k_n \circ f_n$, qui réalise le morphisme canonique
$\Gamma \to \mathrm{R}\Gamma$ lorsque $A^*$ est une résolution de $M_*$ et que
$k_0$ induit l'identité sur $M$.

Si de plus a) $A^i_n = 0$ pour $i > n$, b) $H^n(A^*_n) = 0$ pour $n \geq 1$,
c) $H^{n-1}(A^*_n) = 0$ pour $n \geq 2$, alors un système de Stokes est
déterminé par $k_0$, et $k_0$ en provient si et seulement s'il annule
$\partial_1(\operatorname{Ker}(A^0_1 \xrightarrow{d} A^1_1))$ — ce qui est
automatique si c') $H^0(A^*_1) = \varepsilon^*(A^0_0)$. Il y a alors un
\emph{unique} morphisme $A^*_* \to C^*_{*M}$ prolongeant $k_0$, et, grâce à a),
il est à valeurs dans les cochaînes normalisées :
\[
A^*_* \longrightarrow C^{*!}_{*M} \simeq \Lambda^{*+1}\Phi_* \otimes_k M .
\]

\emph{Exemple.} Pour un anneau à puissances divisées $(S, J)$ et $t \in J$,
l'algèbre $\bigl(\Gamma^*\Phi_{*S} \otimes_S \Lambda^*\Psi_{*S}\bigr)/(T -
t)_{\mathrm{pd}}$ est une résolution de $S$ avec $A^i_j = 0$ pour $i > j$ : on a
donc un système d'intégrales de Stokes à valeurs dans $S$, et un morphisme
canonique vers $\Lambda^{*+1}\Phi_{*S}$, « dont il faudrait étudier les
propriétés multiplicatives ». Le calcul explicite de ces intégrales est fait
pages 136 à 140.

\pagerange{69}{71}
\section*{Comparaison avec le De Rham classique (pages 69 et 71)}

Les données sont : (1) un complexe augmenté $k_X \to A^*_X$ avec intégrales de
Stokes $K_n$ normalisées ; (2) un espace $X$ et une résolution différentielle
$k_X \to \Omega^*$ (par exemple le De Rham d'une variété $C^\infty$, $k =
\mathbb{R}$) ; (3) un sous-ensemble $S(X)'$ du complexe singulier qui donne la
même homologie (les simplexes $C^\infty$) ; (4) pour chaque simplexe $\sigma$,
un « tiré en arrière » des germes de formes au voisinage du support de $\sigma$
vers $A^*_n$, compatible aux opérations simpliciales\footnote{« Supp » est une
lecture conjecturale d'une abréviation qui revient sous toutes les limites
inductives de la page 69.}. Les $K_*$ et (4) définissent un morphisme $\Omega^*
\to S'^*_X$ compatible aux augmentations, d'où un cube commutatif dont la base
est faite d'isomorphismes (image par $\mathrm{R}\Gamma$ d'un carré d'isomorphismes
dans une catégorie dérivée), la flèche verticale de droite étant un
isomorphisme par Cartan, $S^*_X$ étant « homotopiquement fin » ; la face du
haut en hérite. La page 71 est en partie barrée, et le raisonnement n'est
donné qu'à grands traits.

\pagerange{73}{77}
\section*{Le complexe comme noyau de la multiplication par $dT$ (pages 73 à 77)}

Sur $\mathrm{sSet}$, $\Phi_* = i_{0*}(k)$ est un $k$-module muni de la section
$T$, qui définit une injection $k_* \hookrightarrow \Phi_*$. L'algèbre bigraduée
différentielle à puissances divisées
$\mathcal{D}^{**}_* = \Gamma^*_k\Phi_* \otimes_k \Lambda^*\Phi_*$ (différentielle
de bidegré $(-1,+1)$, alternée pour le degré total) est une résolution de $k$ :
$\mathcal{D}^{00}_* = k_*$, et en degré total $n > 0$ la suite $0 \to
\Gamma^n\Phi_* \to \Gamma^{n-1}\Phi_* \otimes \Lambda^1\Phi_* \to \cdots \to
\Lambda^n\Phi_* \to 0$ est exacte, à termes flasques. Soit $\overline{T} = dT
\in \Lambda^1\Phi_*$ ; la multiplication $\overline{T} \wedge$ est de carré nul,
de bidegré $(0,1)$, anticommute à $d$, et a noyau égal à son image. On pose
\[
\mathrm{DR}^{p,q}_* = \operatorname{Ker}\bigl(\mathcal{D}^{p,q+1}_* \xrightarrow{\ \overline{T}\wedge\ } \mathcal{D}^{p,q+2}_*\bigr) \simeq \Gamma^p\Phi_* \otimes \Lambda^q \Psi_* = C^{p,q}_* ,
\]
car le noyau de $\overline{T} \wedge$ sur $\Lambda^{q+1}\Phi_*$ est
$\overline{T} \wedge \Lambda^q\Phi_* \simeq \Lambda^q\Psi_*$. On a les suites
exactes $0 \to \mathrm{DR}^{p,q}_* \xrightarrow{i} \mathcal{D}^{p,q+1}_*
\xrightarrow{u} \mathrm{DR}^{p,q+1}_* \to 0$, où $u = \overline{T}\wedge$.

Sur un $X_*$, avec $\Gamma = \operatorname{Hom}(X_*, -)$ :

(A) $\mathcal{D}^{00}(X_*) \simeq k^{\pi_0(X_*)}$ et
$H^{p,q}(\mathcal{D}(X_*)) = 0$ pour $(p,q) \neq (0,0)$\footnote{Écrit
« si $pq \neq 0$ » ; c'est $(p,q) \neq (0,0)$ que donne la flasquitude des
termes de degré total $> 0$.} ;

(B) les lignes $0 \to \mathcal{D}^{p,0} \xrightarrow{\overline{T}\wedge}
\mathcal{D}^{p,1} \to \cdots$ sont exactes pour $p > 0$, tandis que pour $p =
0$, $\Lambda^\bullet\Phi_*$ muni de $\overline{T}\wedge$ étant une résolution
flasque décalée de $k_*$, la cohomologie est nulle en degrés $\leq 1$ et vaut
$H^{n-1}(X_*, k)$ en degré $n \geq 2$.

La structure multiplicative à puissances divisées de $\mathrm{DR}^{**}$ est
déterminée par la condition que $u : \mathcal{D}^{**} \to \mathrm{DR}^{**}$ soit
compatible à toutes les structures, avec $i(x\,u(y)) = i(x)\,y$, $i(u(y)\,x) =
y\,i(x)$ et $i(x)\,i(y) = 0$ — à la détermination près (E) des produits
$\mathrm{DR}^{0,p} \times \mathrm{DR}^{0,q} \to \mathrm{DR}^{0,p+q}$ ($p,q \geq
1$) et des puissances divisées sur $\mathrm{DR}^{0,p}$, qui sont les seules
données supplémentaires. L'argument (D) : les puissances divisées de
$\mathrm{DR}^{**+} = \operatorname{Ker}(\overline{T}\wedge)$ sont nulles en
ordre $\geq 2$, car localement $y = \overline{T}z$ et $y^{(i)} =
\overline{T}^i z^{(i)}$.

\emph{Problème} (page 77). Trouver des scindages \emph{canoniques} des suites
exactes $0 \to \Gamma^p\Phi_* \otimes \Lambda^{q-1}\Psi_* \to \Gamma^p\Phi_*
\otimes \Lambda^q\Phi_* \to \Gamma^p\Phi_* \otimes \Lambda^q\Psi_* \to 0$,
$p \geq 1$. Un scindage existe, $\Gamma^p\Phi_* \otimes \Lambda^q\Psi_*$ étant
projectif pour $p \geq 1$, dit-il. Pour $p = q = 1$ il est \emph{unique}, car
$\operatorname{Hom}(\Phi_* \otimes \Psi_*, \Phi_*) \simeq
\operatorname{Hom}_k((\Phi_*\otimes\Psi_*)_0, k) = 0$, $\Psi$ étant nul en
degré simplicial $0$ ($\Phi_* = i_{0*}k$ est adjoint à droite de l'évaluation
en $[0]$). La page s'arrête sur « Soit $s : \Phi_* \otimes \Psi_* \to \Phi_*
\otimes \overline{\Phi}_*$ », sans le définir.

\pagerange{79}{90}
\section*{Faisceaux tests sur un topos ; la condition de De Rham (pages 79 à 90)}

\subsection*{Le bicomplexe d'une immersion}

Soient $X$ un topos et $u : \mathbb{Z} \to \Phi$ une immersion telle que a)
$\Psi = \operatorname{Coker} u$ soit plat, b) les $\Lambda^i\Phi$ et $\Gamma^i
\Phi$, $i > 0$, soient flasques et le restent après tout produit tensoriel. On
en déduit que les $\Gamma^i\Phi \otimes \Lambda^j\Phi$ ($i + j > 0$), puis les
$\Gamma^i\Phi \otimes \Lambda^j\Psi$ ($i > 0$), sont flasques. Le bicomplexe
$K^{\bullet\bullet}_u = \Gamma^*\Phi \otimes \Lambda^*\Psi$, avec la
différentielle compatible aux puissances divisées qui est nulle sur
$\Lambda\Psi$ et égale à $\Phi \to \Psi$ sur $\Phi$, est une résolution de
$\mathbb{Z}\{T\}$, où $T \mapsto 1 \in \operatorname{Ker}(\Phi \to \Psi)$ ; et
pour tout Anneau commutatif $k$ sur $X$, $K_{u,k} = \Gamma^*_k\Phi_k \otimes
\Lambda^*_k\Psi_k$ est une résolution de $k\{T\}$, à composantes flasques pour
$i > 0$.

\textbf{Théorème.} \emph{Soit $K^{\bullet\bullet}(X, u, k) = \Gamma(X,
K_{u,k})$. L'homomorphisme canonique $H^{**}(X, u, k) \to H^*(X, k) \otimes_k
k\{T\}$ est un isomorphisme sur la partie de bidegré (complémentaire $p \geq
0$, extérieur $q$), où il s'écrit $H^{p,q}(X,u,k) \simeq H^q(X,k)$, comme au
cas simplicial.}\footnote{La page indexe par (total $i$, extérieur $j$) : isomorphisme
pour $j \leq i$, et « $H^{i,j} = 0$ (et pas un iso) pour $j > i$ ». Comme le
degré extérieur ne dépasse jamais le degré total, ce second membre de phrase
reste obscur ; le dernier facteur est écrit $\mathbb{Z}\{T\}$, ce qui revient
au même puisque $H^*(X,k) \otimes_{\mathbb{Z}} \mathbb{Z}\{T\} \simeq H^*(X,k)
\otimes_k k\{T\}$.} \emph{Plus généralement, pour un morphisme de topos $f : X
\to Y$, on a $f_* K_{u,k} \to \mathrm{R}f_*(k) \otimes_k k\{T\}$, isomorphisme
en cohomologie dans le même domaine.}

\subsection*{Fonctorialités}

$K_u$ est fonctoriel en $k$ et en $u$ : un $\lambda : \Phi \to \Phi'$ tel que
$\lambda u = u'$ induit un morphisme compatible aux structures, isomorphisme
dans la catégorie dérivée, qui le reste après $\Gamma(X,-)$ ou $f_*$. Deux tels
morphismes $\lambda$, $\mu = \lambda + \Theta$ ($\Theta u = 0$) : que peut-on en
dire ? La question est posée page 86 et reste sans réponse.

Les structures de $K$ (page 82) : anneau bigradué alterné pour le degré
extérieur, à bidegrés $\geq 0$, puissances divisées sur l'idéal des bidegrés
non nuls, différentielle compatible, $d : K^{1,0} = \Phi \to K^{1,1} = \Psi$
épimorphique, $K \simeq \Gamma^*\Phi \otimes \Lambda^*\Psi$ libre,
$\Phi$ et $\Psi$ plats ; si $L = \operatorname{Ker} d$, $K$ est une résolution
de $\Gamma^*L$, et l'on suppose le plus souvent $L$ localement libre de rang
$1$. La donnée de $K$ équivaut à celle de $(k, L, \Phi, u : L \hookrightarrow
\Phi)$, et un morphisme à celle d'un morphisme d'anneaux $k \to k'$ et d'un
$\lambda_\Phi$ semi-linéaire respectant les images de $L$.

\subsection*{Indépendance du choix de $u$}

\textbf{Lemme.} \emph{Pour tout $k$-module $N$, $u_N = u \otimes \mathrm{id}_N
: N \to \Phi \otimes_k N$ est un monomorphisme « universel », satisfaisant les
mêmes conditions.} Deux immersions $u$, $u'$ se plongent dans la somme
amalgamée $\Phi \amalg_k \Phi'$, puis dans un $N$ du lemme ; d'où des
isomorphismes $\mathrm{DR}(u) \simeq \mathrm{DR}(N) \simeq \mathrm{DR}(u')$ dans
la catégorie $\mathbb{DR}^{**}$ des anneaux bigradués alternés différentiels à
puissances divisées avec élément de bidegré $(1,0)$, indépendants des choix et
transitifs (par le même argument appliqué à deux choix de $N$). L'argument
vaut pour $f_*\mathcal{DR}^{**}_{\mathrm{pd}}(k,u)$.

\subsection*{La condition de De Rham}

Un Anneau $k$ d'un topos satisfait la \emph{condition de De Rham} s'il existe un
$u : k \to \Phi$ comme ci-dessus ; elle passe à toute $k$-Algèbre, et « le plus
beau » serait que $\mathbb{Z}_X$ y satisfasse. C'est vrai pour le topos
$\widehat{\Delta}$ et pour le topos d'un espace paracompact\footnote{Le second
cas porte « paracompact » souligné comme incertain.} ; pour le topos des
ensembles cosimpliciaux, la réponse inscrite est « Non », soulignée trois fois.
Si $A^{**}$ est une algèbre bigraduée de ce type avec $A^{00} = k$ satisfaisant
la condition, l'augmentation $k\{T\} \to \mathcal{DR}(k,u)$ définit un morphisme
$A^{**} \to \mathcal{DR}(k,u) \otimes_{k\{T\}} A^{**}$, qu'il affirme être un
quasi-isomorphisme.

\pagerange{93}{107}
\section*{Les algèbres de De Rham spéciales (pages 93 à 107)}

\subsection*{Définitions}

Une \emph{structure de De Rham} (à puissances divisées, sous-entendu) est un
anneau bigradué $A^{ij}$, $0 \leq j \leq i$ ($i$ total, $j$ extérieur,
$i - j$ complémentaire), alterné pour le degré extérieur, avec une
différentielle de bidegré $(0,1)$ — une dérivation graduée pour le degré
extérieur —, des puissances divisées sur $A^+ = \sum_{i+j>0} A^{ij}$ compatibles
aux deux graduations ($\gamma^d : A^{ij} \to A^{di,dj}$) et à la différentielle,
et un élément marqué $T \in Z^{10}(A)$. Elle est \emph{spéciale} si $k\{T\}
\xrightarrow{\sim} H^{*0}(A)$, $H^{*j}(A) = 0$ pour $j > 0$, et les $A^{ij}$ sont
$k$-plats\footnote{La condition « $H^{*i} = 0$ » est écrite sans préciser $i$ ;
« $i > 0$ » est la lecture du contexte. Une ligne en interligne, lecture
incertaine, définit encore les structures \emph{flasques} : les $A^{ij}$ ($i
\geq 1$) flasques.}. D'où les catégories $\mathbb{DR}\mathrm{sp}(\mathfrak{X},
k)_0$ (morphismes respectant toutes les structures),
$\mathbb{DR}\mathrm{spfl}$, $\mathbb{DR}\mathrm{spufl}$ (flasques,
universellement flasques), et leurs catégories de fractions obtenues en
inversant les quasi-isomorphismes.

\subsection*{Réduction à l'algèbre standard}

Pour $A$ spéciale, $T$ définit $k \to A^{10} = \Phi$, d'où une suite exacte $0
\to k \to \Phi \to \Psi \to 0$ avec $\Psi$ plat, et un morphisme
$\mathrm{DR}(\Phi, T) = \Gamma^*\Phi \otimes \Lambda^*\Psi \to A$. Tout
morphisme $A \to A'$ s'inscrit dans un carré avec $\mathrm{DR}(\Phi,T) \to
\mathrm{DR}(\Phi', T')$.

\textbf{Proposition} (page 97). \emph{La catégorie localisée de
$\mathbb{DR}\mathrm{sp}(\mathfrak{X}, k)_0$ est équivalente à la catégorie
ponctuelle ; de même pour $\mathbb{DR}\mathrm{spufl}(\mathfrak{X},k)_0$, si les
$\Gamma^i\Phi$ ($i > 0$) d'un $\Phi$ plat universellement flasque sont encore
universellement flasques.} Elle se ramène à la même assertion pour la catégorie
des extensions $0 \to k \to \Phi \to \Psi \to 0$ à $\Psi$ plat, que la page dit
« on sait » équivalente à la catégorie ponctuelle\footnote{Ce dernier point
n'est pas démontré dans le dossier ; c'est la même question que celle des deux
morphismes $\lambda$, $\mu$ de la page 86 et de la simple connexité de la page
120. Pour $\widehat{\Delta}$, la page 99 justifie l'hypothèse par
« flasque $+$ plat $\Rightarrow$ homotope à $0$ $\Rightarrow$ universellement
flasque et les $\Gamma^i$ universellement flasques », que la marge déclare
elle-même « faux tel quel, à arranger ! ».}.

\subsection*{Sur $\widehat{\Delta}$}

Soit $i_n : P \to \widehat{\Delta}$ le point $K \mapsto K_{[n]}$.
\textbf{Lemme.} \emph{$i_{0*}(M) \simeq \Phi_* \otimes_k M$, avec $\Phi_*(I) =
k^I$.} Car $i_{0*}(M)(I) = M^I$. \textbf{Corollaire.}
$\operatorname{Hom}(A_*, \Phi_* \otimes_k M) \simeq \operatorname{Hom}_k(A_0, M)$,
d'où un morphisme canonique $A_* \to \Phi_* \otimes A_0 = i_{0*}i_0^*A_*$.
Appliqué à $k_* \hookrightarrow A^{1,0}_*$, cela relie toute algèbre de De Rham
spéciale, « en trois crans », à l'algèbre standard $\mathrm{DR}(\Phi_{*k}, T)$,
« qui semble la plus économique de toutes » (page 101) :
\[
A^{**}_* \longleftarrow \mathrm{DR}(A^{10}_*, T) \longrightarrow \mathrm{DR}(A^{10}_0 \otimes \Phi_*, T) \longleftarrow \mathrm{DR}(\Phi_{*k}, T) ,
\]
diagramme qui reste dans les flasques si $A^{10}_*$ l'est.

Il introduit ensuite des catégories $\mathbb{DR}\mathrm{sp}'$,
$\mathbb{DR}\mathrm{spfl}'$ ayant plus de morphismes : les homomorphismes
d'anneaux à puissances divisées entre complétés $\widehat{A} = \prod A^{ij}$
respectant $T$. Les groupoïdes engendrés seraient « triviaux » (réunion
d'objets reliés par des isomorphismes transitifs), et équivalents « au
groupoïde connexe défini par deux groupes, que nous nous proposons de
déterminer » dans le cas de $\widehat{\Delta}$ (page 105). Les pages 105 et 107
commencent la détermination des morphismes de
$\widehat{\mathrm{DR}}_{*k}$ dans une algèbre de De Rham : des composantes
$v(p,q)$, avec $v(1,0)(T) = T$, et, page 107 — au crayon très pâle, lue par
fragments seulement —, un élément $\xi_0$ en degré simplicial $1$ et des
relations de commutation entre ses images par les deux dégénérescences
$\Delta_2 \to \Delta_1$. Le premier groupe est déterminé aux pages 161 à
179 ; le second n'apparaît pas dans le dossier\footnote{Le feuillet de la page
104 ne porte qu'un dessin du diagramme ci-dessus au pied d'une page
dactylographiée étrangère ; les pages 102, 106, 108 à 115 et 117 à 119 sont
des versos dactylographiés, une feuille blanche, le début d'une lettre
personnelle (page 117) et une lettre d'un tiers (page 119), qui ne sont pas
reproduits.}.

\pagerange{110}{120}
\section*{Les théories de De Rham admissibles (pages 110 à 120)}

On considère les foncteurs $C^{**} : \mathrm{sSet}^\circ \to \mathfrak{D}$, où
$\mathfrak{D}$ est la catégorie des $S_k$-algèbres bigraduées alternées
différentielles augmentées à puissances divisées, qui transforment limites
inductives en limites projectives ; ils sont représentables, c'est-à-dire donnés
par des objets simpliciaux $C^{**}_*$ de $\mathfrak{D}$. À un tel $C$ est
associé le foncteur cohomologique $\mathcal{H}^{**}_C(X) = H^{**}(C^{**}(X))$.
$C$ est \emph{admissible} si, pour tout $X$, l'image de $H^{**}$ modulo les
modules négligeables (le foncteur $\tau_\omega$ des pages 253 à 258) est dans
l'image essentielle des $k$-modules, $\mathcal{H}_C(X) \approx H^*_C(X)
\otimes_k S_k$.

Précisément : dans le diagramme $\mathrm{R}\Gamma(X, S_k) \to \mathrm{R}\Gamma(X,
C^{**}) \leftarrow C^{**}(X)$, on veut que les deux flèches deviennent des
isomorphismes après $\tau_\omega$, ce qui entraîne que les morphismes induits
en cohomologie sont à noyau et conoyau négligeables. Pour $X = \Delta_n$ cela
se réduit à :
3a) $S^i_k \to H^{i,0}(C_n)$ est un isomorphisme pour $i \geq \lambda(n)$ ;
3b) $H^{i,p}(C_n) = 0$ pour $i \geq \lambda(p)$, $p > 0$\footnote{Ici $i$ est le
degré total et $p$ le degré extérieur.} — avec $\lambda$ strictement
croissante, dont « on suspecte » qu'on peut la prendre indépendante de $n$.
Dans ce cas la flèche horizontale devient un isomorphisme après $\tau_\omega$,
chaque $C^{i*}_*$ étant une résolution de $(S^i_k)_*$ jusqu'en degré $p$ pour
$i$ grand. Pour la flèche verticale suffit (4) : si $\pi(i)$ est le plus petit
entier tel que $C^{ip}_*$ soit flasque pour $p \geq \pi(i)$, on veut $\pi(i) \to
+\infty$.

Il aimerait un théorème d'unicité. Les morphismes entre théories admissibles
sont déjà des quasi-isomorphismes modulo négligeables ; on passe au groupoïde
engendré, et l'on « aimerait prouver que ce groupoïde est connexe et simplement
connexe » : a) deux théories sont dominées par une troisième ; b) deux
morphismes parallèles sont égalisés par un morphisme. Un grand « ? » en marge,
et : « il y aurait plus de chances avec catégorie de fractions à droite ??? ».

\pagerange{122}{134}
\section*{La version duale : coalgèbres et chaînes de De Rham (pages 122 à 134)}

\subsection*{Schémas affines en algèbres}

Une nouvelle rédaction, intitulée « Complexes duaux de De Rham à
p.d. »\footnote{Le mot « duaux » est une lecture incertaine.}, part d'une
$k$-coalgèbre $s$ associative, unitaire, commutative, d'une sous-coalgèbre
$\bar s$, du quotient $s^+ = s/\bar s$, d'un système de « copuissances
divisées » $\gamma^*_i : s^+ \to \operatorname{Sym}^i_k(s^+)$, et d'une forme
$t : s^+ \to k$. Le foncteur $V_k(s) : k' \mapsto s^{k'} =
\operatorname{Hom}_k(s, k')$ est à valeurs dans les anneaux commutatifs ;
$V_k(s^+)$ en est un Idéal, muni de puissances divisées par les
$\gamma^*_i$, et $t$ en est une section. On obtient un objet simplicial dans les
schémas affines en algèbres différentielles graduées alternées à idéal à
puissances divisées :
\[
\mathcal{A}_n(s,s^+,\gamma,t)(k') = s^{k'}\{\{X_0,\ldots,X_n\}\}[dX_0,\ldots,dX_n] \big/ \Bigl(\textstyle\sum X_i - t,\ \sum dX_i\Bigr)_{\mathrm{pd}} ,
\]
où les doubles accolades désignent les séries formelles à puissances
divisées. Comme foncteur en $k'$ c'est $\operatorname{Hom}_k(s,k')^{\mathbb{N}^n
\times \mathfrak{P}([0,n-1])}$, donc
\[
\mathcal{A}_n = V\bigl(\mathcal{A}^n\bigr), \qquad \mathcal{A}^n = s \otimes_k k[Y_0,\ldots,Y_{n-1}] \otimes_k \Lambda^*(y_0,\ldots,y_{n-1}) .
\]
a) Chaque $\mathcal{A}^n$ est une $k$-coalgèbre graduée alternée
codifférentielle (de degré $-1$), avec $\bar s \hookrightarrow \mathcal{A}^n$ et
des copuissances divisées ; b) les $\mathcal{A}^n$ forment un objet
cosimplicial ; c) $\mathcal{A}_* = V_k(\mathcal{A}^*)$.

\subsection*{Chaînes de De Rham}

d) Pour un ensemble simplicial $\mathcal{X}$,
\[
\operatorname{Hom}\bigl(\mathcal{X}, \mathcal{A}_*(k')\bigr) \simeq \varprojlim_{\sigma \in \Delta/\mathcal{X}} \mathcal{A}_{\dim\sigma}(k') \simeq \operatorname{Hom}_k\Bigl(\varinjlim_{\sigma \in \Delta/\mathcal{X}} \mathcal{A}^{\dim\sigma}, k'\Bigr) ,
\]
et l'on pose $C^{\mathrm{DR}}_*(\mathcal{X}; s, s^+, \gamma, t) =
\varinjlim_{\sigma} \mathcal{A}^{\dim\sigma}$ : une $k$-coalgèbre graduée
alternée codifférentielle avec un morphisme $\bar s \to C^{\mathrm{DR}}_*$ (non
nécessairement injectif : il est nul si $\mathcal{X} = \varnothing$) et des
copuissances divisées. Elle commute aux limites inductives quelconques en
$\mathcal{X}$, et les cochaînes s'en déduisent par
$C^*_{\mathrm{DR}}(\mathcal{X}; s^{k'},\ldots) = \operatorname{Hom}_k(C^{\mathrm{DR}}_*,
k')$. e) Elle est covariante en $(s, s^+, \gamma, t)$ et commute aux limites
inductives filtrantes et à l'extension de la base. f) Si $s$ est graduée et
$t$ homogène de degré $\delta$, tout devient bigradué ($\deg X_r = \deg dX_r =
\delta$).

g) On prend $s = k[u]$, $\bar s = k$, d'où $s^{k'} = k'\{\{T\}\}$ ; les
copuissances divisées sont celles qui correspondent aux puissances divisées de
$k'\{\{T\}\}^+$, et $t = T$, c'est-à-dire $t(u^n) = 1$ si $n = 1$, $0$ sinon. On
obtient $\mathcal{A}^*_{k*}$ avec $\operatorname{Hom}_k(\mathcal{A}^*_k, k') \simeq
\mathcal{A}_*(k')^\wedge$ — le complété du complexe des pages 5 à 8 —, et la
coalgèbre bigraduée $C^{\mathrm{DR,pd}}_*(\mathcal{X}, k)$ qui redonne par
dualité toutes les structures de $C^*_{\mathrm{DR,pd}}(\mathcal{X}, k')$. Enfin
tout se déduit du cas $\mathbb{Z}$ :
\[
C^{\mathrm{DR,pd}}_*(\mathcal{X}, k) \simeq C^{\mathrm{DR,pd}}_*(\mathcal{X}) \otimes_{\mathbb{Z}} k , \qquad C^*_{\mathrm{DR,pd}}(\mathcal{X}; M) = \operatorname{Hom}\bigl(C^{\mathrm{DR,pd}}_*(\mathcal{X}), M\bigr)
\]
pour tout groupe abélien $M$. C'est l'objet « le plus fin » annoncé page 14.

\pagerange{136}{140}
\section*{Dualisation des homomorphismes de Stokes (pages 136 à 140)}

Les intégrales de Stokes $\pi_n : \mathcal{A}^n_n \to s^{k'}$ se dualisent en
$\kappa_n : s \to (\mathcal{A}^*)_n$ en degré $n$, avec $\kappa_0 = \mathrm{id}$
et la condition de Stokes $d_{\mathcal{A}^*}\kappa_n = \partial\,\kappa_{n-1}$,
d'où des morphismes $C^{\mathrm{sing}}_*(\mathcal{X}, s) \to
C^{\mathrm{DR}}_*(\mathcal{X}; s, s^+, \gamma, t)$\footnote{Le haut de la page
136 est biffé et semble appartenir à un autre fil (faisceaux sur un ensemble
simplicial). Dans le carré de Stokes la flèche de gauche est notée
$\pi_{n-1}$, là où l'on attend $\kappa_{n-1}$.}.

La dualité est explicite (page 138) : en degré cotangent $0$,
$k\{\{T\}\}\{\{X_0,\ldots,X_n\}\}/(\sum X_i - T)_{\mathrm{pd}} \simeq
k\{\{X_0,\ldots,X_n\}\}$ est dual de $k[Z_0,\ldots,Z_n] \hookrightarrow
k[U, Y_0,\ldots,Y_n]$, $Z_i \mapsto U + Y_i$, par l'accouplement
$\langle X^{(i)}, Z^j\rangle = \delta_{i,j}$ ; l'image est l'intersection des
noyaux des opérateurs différentiels divisés $\bigl(\sum_i \partial/\partial Y_i
- \partial/\partial U\bigr)^{(m)}$. En degré total égal au degré cotangent,
$\Lambda^*[dX_0,\ldots,dX_n]/(\sum dX_i)$ est dual du sous-espace des vecteurs de
somme nulle, noyau du produit intérieur par la forme « somme des
coordonnées ».

L'intégrale elle-même est calculée page 140 : sur $(\mathcal{A}^*_n)^n \simeq
S\{X_0,\ldots,X_{n-1}\}\, dX_0\wedge\cdots\wedge dX_{n-1}$,
\[
\kappa^n\bigl(X_0^{(i_0)} \cdots X_{n-1}^{(i_{n-1})}\, dX_0 \wedge \cdots \wedge dX_{n-1}\bigr) = \gamma_{i_0 \ldots i_{n-1}}\; t^{(n + \sum i_\alpha)} ,
\]
avec $\gamma_{i} = \frac{(n+\sum i_\alpha)!}{\prod i_\alpha!} \int_{\sigma_n}
X_0^{i_0}\cdots X_{n-1}^{i_{n-1}}\,dX_0\cdots dX_{n-1}$ sur le simplexe $\sigma_n
= \{X_\alpha \geq 0,\ \sum X_\alpha \leq 1\}$. La page calcule ensuite $\gamma$
par récurrence sur $n$ au moyen de Stokes, et la récurrence s'arrête au bas de
la page. En fait $\gamma_i = 1$ pour tout $i$ : c'est l'intégrale de Dirichlet
$\int_{\sigma_n} \prod X_\alpha^{i_\alpha} = \prod i_\alpha! / (n + \sum
i_\alpha)!$\footnote{Cette évaluation est la nôtre ; elle suppose l'orientation
standard du simplexe, à laquelle le signe est attaché. La page n'écrit, en
accolade, que l'encadrement « $0 \leq \cdot \leq \frac{1}{n!}$,
$\gamma_{0,\ldots,0} \overset{?}{=} \frac{1}{n!}$ », qui est le cas $i = 0$.
Les coefficients $c_{i,l}$ de la récurrence sont ceux du produit des
puissances divisées, ce que la page ne dit pas.}. L'intégrale, sur le simplexe
de taille $t$, d'un monôme en puissances divisées est donc la puissance divisée
de $t$ de l'ordre total :
\[
\int_{t\sigma_n} X_0^{(i_0)} \cdots X_{n-1}^{(i_{n-1})}\, dX_0\cdots dX_{n-1} = t^{(n + \sum i_\alpha)} ,
\]
ce qui explique pourquoi l'intégration de Stokes est définie sans dénominateur.

\pagerange{143}{151}
\section*{Le yoga coalgèbres $\mapsto$ Algèbres (pages 143 à 151)}

On travaille dans le topos $\widehat{C}_k$ des préfaisceaux sur la catégorie
opposée des $k$-algèbres, où le foncteur oubli est une $k$-Algèbre
$\mathcal{O}_k$. Pour un $k$-module $M$, $V(M)(A) = \operatorname{Hom}_k(M, A)$
est un $\mathcal{O}_k$-Module, représenté par $\operatorname{Sym}^*(M)$ ; d'où
$\operatorname{Hom}_{\widehat{C}_k}(V(N), F) \simeq F(\operatorname{Sym}^* N)$ et
\[
\operatorname{Hom}_{\widehat{C}_k}(V(N), V(M)) \simeq \operatorname{Hom}_k(M, \operatorname{Sym}^* N) .
\]
Un morphisme $u^*$ est homogène de degré $n$ si et seulement si le $u$
correspondant tombe dans $\operatorname{Sym}^n N$ (on teste sur $A = k[T]$,
$\lambda = T$) ; les morphismes homogènes de degré $1$ sont donc exactement les
$\operatorname{Hom}_k(M, N)$, nécessairement $\mathcal{O}_k$-linéaires, et $V$
est pleinement fidèle (contravariant). De même, $V$ transforme sommes en
produits, et les accouplements multilinéaires $\prod V(N_\alpha) \to V(M)$
correspondent aux $M \to \bigotimes N_\alpha$ ; une loi bilinéaire sur $V(M)$
est donc une comultiplication $\Delta : M \to M \otimes M$, associative,
commutative, unitaire si la loi l'est ; une unité de $V(M)$ est une counité $M
\to k$, une augmentation $V(M) \to \mathcal{O}_k$ un élément $e$ avec $\Delta e =
e \otimes e$. Une limite inductive quelconque de coalgèbres est une coalgèbre.

Une graduation de type $Z$ de $M$ revient à décomposer $V(M)$ en produit, et,
si $Z$ est un groupe, à une action du groupe diagonalisable $D(Z) =
\operatorname{Spec} k[Z]$. Une coalgèbre graduée correspond à une loi
compatible aux graduations ; avec $Z \to \mathbb{Z}/2$ on définit les
coalgèbres anticommutatives et \emph{alternées} ($x^2 = 0$ en degré impair),
alterné entraînant anticommutatif ; ces conditions passent aux limites. Les
codérivations et coantidérivations correspondent aux dérivations et
antidérivations de $V(M)$ ; un endomorphisme de degré $\nu$ de $M$ donne un
endomorphisme de degré $-\nu$ de $V(M)$, ce qui n'importe pas pour les
dérivations, où seule la parité compte. Les « coproduits divisés » sont à
traiter de même, après le cas ordinaire.

\pagerange{152}{159}
\section*{Coalgèbres et comodules (pages 152 à 159)}

Si $C$ est un comodule sur la coalgèbre $s$, $V(C)$ est un Module sur $V(s)$,
donc $C$ est un module ordinaire sur l'algèbre duale $s^k =
\operatorname{Hom}_k(s, k)$, $u \in s^k$ opérant par $C \xrightarrow{\Delta_C}
C \otimes s \xrightarrow{\mathrm{id} \otimes u} C$. On récupère la structure de
comodule si $C \otimes s \to \prod_{u \in s^k} C$ est injectif, ce qui est le cas
si $s$ est \emph{libre} sur $k$ (sans hypothèse de platitude sur $C$).
Exemple : $s = k[U]$, $s^k = k\{\{T\}\}$ ; $C = k[U, Y_0,\ldots,Y_n]$ est un
module sur $C^k = k\{\{T\}\}\{\{X_0,\ldots,X_n\}\}$, $T^{(i)}X^{(i_0)}_0\cdots$
opérant par les opérateurs différentiels divisés $D_U^{(i)} D_{Y_0}^{(i_0)}
\cdots$, localement nilpotents.

Pour $s$ cocommutative, unitaire et libre de base $(e_\alpha)$, les modules qui
proviennent d'un comodule sont ceux pour lesquels $u \mapsto u.x$ est continu
pour la topologie de la convergence simple sur $s^k \simeq k^A$ et la topologie
discrète sur $C$ ; la catégorie des $s$-comodules gradués équivaut à celle des
$s^k$-modules gradués « localement nilpotents » en ce sens\footnote{Ce sont
les modules qu'on appelle aujourd'hui \emph{rationnels} (ou discrets) sur
l'algèbre duale d'une coalgèbre.}. Si $u \in (s^k)_n$, $u_C$ \emph{diminue} les
degrés de $n$ ; pour qu'il les augmente il faut changer le signe des degrés,
$V(C)^m = V(C_{-m})$ — ce que reprennent les notes de la page 159.

\pagerange{161}{171}
\section*{Les homomorphismes de l'algèbre simpliciale $\mathcal{A}_*$ (pages 161 à 171)}

Notons $\mathcal{A}_S = \mathcal{A}(S, S^+, t)_*$ l'algèbre simpliciale des
pages 5 à 8, pour un anneau $S$ à idéal à puissances divisées $S^+$ et $t \in
S^+$ ; on a $(\mathcal{A}_S)_0 = S$ et $(\mathcal{A}_S)_1 \simeq
S\{X_0\}[dX_0]$, avec les deux augmentations $\alpha(X_0) = 0$ et
$\beta(X_0) = t$ (faces vers $(\mathcal{A}_S)_0$).

\textbf{Théorème} (page 161). \emph{Les homomorphismes d'anneaux simpliciaux
$\varphi_* : \mathcal{A}_S \to \mathcal{A}_{S'}$ — sans respecter a priori ni
graduation, ni filtration, ni puissances divisées — correspondent
biunivoquement aux systèmes $(\varphi_0, A'_1, (B'_i)_{i \geq 1})$ formés d'un
homomorphisme d'anneaux $\varphi_0 : S \to S'$, de $A'_1 = \sum_{j \geq 1}
a'_j X_0^{(j)} \in \operatorname{Ker}(\alpha)$ et de $B'_i = \sum_{j \geq 0}
b'_{ij}X_0^{(j)} \in S'\{X_0\}$, soumis aux conditions :}
\emph{1) $\varphi_0(t^{(i)}) = \varphi_0(t)^{(i)}$ — $\varphi_0$ commute
nécessairement aux puissances divisées sur $(t)_{\mathrm{pd}}$ ;}
\emph{2) $\beta(A'_1) = \varphi_0(t)$, i.e.\ $\sum_j a'_j\, t'^{(j)} =
\varphi_0(t)$ ;}
\emph{3) $(2B'_i)\,J' = 0$ pour tout $i$, où $J' \subset S'$ est l'idéal
engendré par les coefficients $a'_j$ et $b'_{jk}$.}

On reconstitue $\varphi_1$ par $\varphi_1|_S = \varphi_0$,
$\varphi_1(X_0^{(i)}) = A_1'^{(i)} + B'_i\,dX_0$, $\varphi_1(dX_0) = dA'_1$, et
$\varphi_n$ en écrivant $(\mathcal{A}_S)_n = S\{U_0,\ldots,U_{n-1}\}[dU_0,
\ldots,dU_{n-1}]$, $U_r = X_0 + \cdots + X_r$, et en posant $\varphi_n(U_r^{(i)})
= A_1'^{(i)}\{U_r\} + B'_i\{U_r\}\,dU_r$. La condition 3) ne sert qu'à partir
de $n = 2$, pour que les images de variables différentes commutent\footnote{La
page ne fait que l'affirmer ; voici la raison. Si $a + b$ et $a' + b'$ sont deux
éléments dont les parties de degré $1$ sont $b = B_r\,dU_r$ et $b' = B'_s\,dU_s$
($r \neq s$), leur commutateur est $bb' - b'b = 2\,bb' = 2B_rB'_s\,dU_r\,dU_s$ ;
de même avec $\varphi_n(dU_s) = A_1'^{\mathrm{der}}\{U_s\}\,dU_s$. D'où
l'annulation de $2B'_i$ contre les coefficients des $A'$ et des $B'$.}.
Si $\varphi_0(t)$ est régulier, ou si le coefficient $a'_1$ est inversible, 3)
signifie $2B'_i = 0$.

\emph{Corollaire 1.} Si $2$ est inversible dans $S'$ (sous l'une de ces
conditions), $\varphi_*$ respecte la graduation, et si $S^+ =
(t)_{\mathrm{pd}}$, les puissances divisées.
\emph{Corollaire 2.} $\varphi_*$ commute aux puissances divisées si et
seulement si $B'_i = A_1'^{(i-1)}B'_1$ ; il respecte les graduations si et
seulement si $B'_i = 0$ (et alors il commute aux puissances divisées).
\emph{Corollaire 3.} Il respecte la filtration croissante grossière si et
seulement si $A'_1 = a'_1X_0$ et $B'_i \in \mathrm{Fil}^{i-1}$ ; graduations et
filtrations ensemble : $A'_1 = a'_1X_0$ et $B'_i = 0$, et ces $\varphi$ sont
les prolongements canoniques des $(\varphi_0, \lambda)$ avec
$\lambda t' = \varphi_0(t)$\footnote{La page écrit « $B'_i = 0$, $i \geq 0$ » ;
les $B'_i$ ne sont définis que pour $i \geq 1$.}.
\emph{Corollaire 4.} $\varphi_*$ respecte toujours la filtration décroissante
associée au degré extérieur.

\emph{Le cas $S = k\{T\}$, $S' = k'\{T\}$.} $\varphi_0$ est donné par
$\varphi_{00} : k \to k'\{T\}$, $\lambda \mapsto \sum_i \delta_i(\lambda)
T^{(i)}$, avec $\delta_0$ homomorphisme, $\delta_i(1) = 0$ pour $i \geq 1$ et
$\delta_i(\lambda\lambda') = \sum_{j+l=i} \binom{i}{j} \delta_j(\lambda)
\delta_l(\lambda')$\footnote{La page écrit la règle de Leibniz sans le
coefficient binomial ; il est imposé par $T^{(j)}T^{(l)} = c_{j,l}T^{(j+l)}$.
Les $\delta_i$ forment ce qu'on appelle aujourd'hui une \emph{dérivation de
Hasse–Schmidt} (écrite ici dans la base des puissances divisées).}, et par
$\varphi_0(T)$, fixé par 2). Identifiant $S'\{X_0\} = k'\{T\}\{X_0\} \simeq
k'\{X_0, X_1\}$ ($T = X_0 + X_1$), on a $A'_1 = \sum_{i \geq 1, j \geq 0}
\alpha_{ij}X_0^{(i)}X_1^{(j)}$ et $\varphi_0(T) = \sum_i \alpha_{i,0}T^{(i)}$.

\textbf{Corollaire 5.} \emph{Les homomorphismes $\mathcal{A}(k\{T\}, k\{T\}^+, T)
\to \mathcal{A}(k'\{T\}, k'\{T\}^+, T)$ sont donnés par $\varphi_{00} : k \to
k'\{T\}$, $A'_1 \in k'\{X_0,X_1\}$ sans terme en $X_1^{(j)}$ seul, et des $B'_i
\in k'\{X_0,X_1\}$, avec la seule condition $(2b'_{ij})J' = 0$.} Ils commutent
aux puissances divisées si $B'_i = A_1'^{(i-1)}B'_1$, aux graduations
extérieures si $B'_i = 0$, aux filtrations croissantes grossière ou fine sous
des conditions explicites (pages 169 et 171), au degré total si $\varphi_0(k)
\subset k'$, $A'_1 = \alpha'_0X_0$ et $B'_i$ homogène de degré $i-1$ en
$(X_0,X_1)$. Ceux qui respectent les deux graduations sont exactement les
prolongements des $u_{\varphi_{00}, \alpha'_0} : k\{T\} \to k'\{T\}$, $T \mapsto
\alpha'_0 T$ — « vous voyez un minimum ! »\footnote{Sous e), page 171, la
condition sur $\varphi_0$ et $A'_1$ qui devrait précéder celle des $B'_i$ ne se
lit pas ; seule la ligne des $B'_i$ est écrite.}.

\pagerange{173}{179}
\section*{Le schéma en groupes des automorphismes (pages 173 à 179)}

On complète : $\widehat{\mathcal{A}}_{k,n} \simeq k\{\{T\}\}\{\{X_0,\ldots,
X_{n-1}\}\}[dX_0,\ldots,dX_{n-1}]$ ; les résultats sont les mêmes, $A$ et les
$B_i$ devenant des séries à puissances divisées dans $k\{\{X_0,X_1\}\}$. Pour un
endomorphisme $\varphi_*$ de la $k$-algèbre différentielle simpliciale
$\widehat{\mathcal{A}}_{k*}$, écrivons $A = \sum_{i \geq 1, j \geq 0}
\alpha_{ij}X_0^{(i)}X_1^{(j)}$ et $\varphi_0(T) = \sum_{j \geq 1} \lambda_j
T^{(j)}$ ; on a $\alpha_{1,0} = \lambda_1$, l'\emph{indicateur}
$\Delta(\varphi_*)$ de $\varphi_*$\footnote{Le mot est hésitant sur la page :
« discriminant » biffé, « déterminant » ajouté, « indicateur » doublement
souligné et peut-être lui aussi barré ; c'est celui de l'énoncé.}.

\textbf{Théorème.} \emph{$\varphi_*$ est un automorphisme si et seulement si
son indicateur est inversible. $\Delta$ est alors multiplicatif, $\Delta(1) =
1$, d'où un homomorphisme $\Delta : \Omega(k) =
\operatorname{Aut}(\widehat{\mathcal{A}}_{k*}) \to k^* = \mathbf{G}_m(k)$.}

Les éléments de $\Omega(k)$ sont les systèmes $(a_{jk})_{j \geq 1, k \geq 0}$,
$(b_{ijk})$ d'éléments de $k$ (coefficients de $A$ et des $B_i$ en puissances
divisées de $T$), soumis aux seules conditions $a_{10} \in k^*$ et $2b_{ijk} =
0$ (la condition 3 du théorème, $a_{10}$ étant inversible). Le foncteur
$\Omega$ est donc représentable par le schéma affine
\[
\Omega \simeq \operatorname{Spec}\Bigl(\mathbb{Z}\bigl[(A_{jk}), (B_{ijk})\bigr] \big/ (2B_{ijk})\Bigr)\bigl[A_{10}^{-1}\bigr] ,
\]
c'est-à-dire $\mathbf{G}_m \times \mathbf{G}_a^{(\mathbb{N}^* \times
\mathbb{N}) \smallsetminus \{(1,0)\}} \times {}_2\mathbf{G}_a^{\mathbb{N}^* \times
\mathbb{N} \times \mathbb{N}}$ comme schéma, où ${}_2\mathbf{G}_a$ est le noyau de
la multiplication par $2$\footnote{La page écrit l'indice ${}_2$ aussi devant le
premier facteur $\mathbf{G}_a$, celui des $a_{jk}$, qui ne sont soumis à aucune
condition hors $a_{10}$ ; la transcription signale que c'est peut-être un
trait parasite. Nous suivons la description par les coefficients, qui est sans
ambiguïté.}. C'est un schéma en groupes sur $\mathbb{Z}$ — « par abus de
langage, le schéma en groupes d'automorphismes de l'anneau différentiel
simplicial $\widehat{\mathcal{A}}_{\mathbb{Z}*}$ ». La section $\delta :
\mathbf{G}_m \to \Omega$, $\lambda \mapsto (A = \lambda X_0, B_i = 0)$, soit
$T \mapsto \lambda T$ et $X_i \mapsto \lambda X_i$, fait de $\Omega$ un produit
semi-direct de $\mathbf{G}_m$ et d'un groupe unipotent fibre par fibre.

Si $2$ est inversible dans $k$, les automorphismes respectent le degré
extérieur ; leurs effets sur la cohomologie des $\operatorname{Hom}(X_*,
\widehat{\mathcal{A}}_{k*})$ sont donc des opérations de degré $0$, qui se
factorisent par $\Delta$ : des multiplications par des puissances de $\lambda
\in k^*$ fixées par le bidegré\footnote{La page écrit « mult.\ par $\lambda^i$
dans $H^i$ » au bout d'une phrase en grande partie illisible ; l'exposant exact
dépend de l'indexation retenue, et nous ne le fixons pas.}. Aucune opération
intéressante ne sort donc de là.

\pagerange{181}{199}
\section*{La caractéristique $2$, et les endomorphismes de $\mathcal{A}_1$ (pages 181 à 199)}

« Si donc on veut trouver des opérations cohomologiques intéressantes », il
faut se placer dans le cas où $2$ n'est pas inversible : « on peut espérer
trouver alors une description des carrés de Steenrod mod $2$ ». On peut aussi
appliquer le foncteur $\tau_\omega$, $M \mapsto M_\omega$, qui fait de
$X_* \mapsto \operatorname{Hom}(X_*, \mathcal{A}_*)$ un foncteur à valeurs dans
les algèbres différentielles de la catégorie abélienne $\mathcal{M}_\omega$
(« de type Artin », avec un $\otimes$), et chercher les types d'anneaux et de
foncteurs en oubliant la graduation. Le dossier ne revient pas sur ces
espoirs.

\textbf{Proposition} (page 185). \emph{Soient $S$ un anneau commutatif et $A =
S\{X_0\}[dX_0]$, augmentée par $\alpha$. Les endomorphismes de $S$-algèbre
différentielle augmentée de $A$ sont donnés par $\varphi(X_0^{(i)}) = A_1^{(i)} +
B_i\,dX_0$, $\varphi(dX_0) = A_1'\,dX_0$, où $A_1 \in S\{X_0\}^+$ et les $B_i \in
S\{X_0\}$ sont arbitraires ; $\varphi$ commute aux puissances divisées si et
seulement si $B_i = A_1^{(i-1)}B_1$.} \emph{Corollaire : s'il respecte la
graduation, il respecte les puissances divisées.} Pour $S = k\{T\}$, la
compatibilité à la seconde augmentation $\beta$ ($\beta(X_0) = T$) s'écrit
$\beta(A_1) = T$ ; posant $f = A_1 - X_0$ et identifiant $S\{X_0\} \simeq k\{X_0,
X_1\}$, les deux conditions disent $f \in (X_0)_{\mathrm{pd}} \cap
(X_1)_{\mathrm{pd}}$, soit $f = \sum_{i,j \geq 1} f_{ij}X_0^{(i)}X_1^{(j)}$.

\textbf{Proposition} (pages 189 à 191). \emph{Les endomorphismes de la
$S$-algèbre différentielle $\mathcal{A}_1$ compatibles aux deux augmentations
sont donnés par $f = \sum_{i,j \geq 1} f_{ij}X_0^{(i)}X_1^{(j)}$ et des $B_i \in
k\{X_0,X_1\}$, par $\varphi(X_0^{(i)}) = (X_0 + f)^{(i)} + B_i\,dX_0$ et
$\varphi(dX_0) = dX_0 + df$.}\footnote{La page écrit $\varphi(dX_0) = 1 +
df$ ; c'est $d(X_0 + f)$.} Échanger $X_0$ et $X_1$ change $f$ en $-f$ et les
$B_i$ en $C_i = \sum_{j \leq i} (-1)^{j+1}T^{(i-j)}B_j$.

\textbf{Proposition} (page 193). \emph{Un tel $\varphi$ se prolonge de façon
unique en des $\varphi_n$ compatibles à toutes les $\sigma : \Delta_n \to
\Delta_1$.} Le corollaire — les endomorphismes de l'algèbre simpliciale $A_*$
correspondent aux $\varphi$ — est barré, avec en marge « faux tel quel ; vrai si
$2B_i = 0$ ». C'est exactement la condition 3) du théorème de la page 161, qui
en est la version corrigée. La démonstration (pages 195 à 199) utilise les
surjections $\sigma_r : \Delta_n \to \Delta_1$, $A_*(\sigma_r)(X_0) = U_r = X_0 +
\cdots + X_r$, le lemme $A_n^0 \simeq S\{U_0,\ldots,U_{n-1}\}$, et l'unique
morphisme d'algèbres $U_r^{(i)} \mapsto \sigma_r^*(\xi_i)$ ; il reste à vérifier
que les images des $dU_r$ sont de carré nul, anticommutent entre elles et
commutent aux $\xi_{s,i}$ — c'est là que surgit, dans la note marginale,
l'annulation de $2b_{ik}$. Le lemme final sur les morphismes d'algèbres
différentielles $S\{U\}[dU] \to B$ s'interrompt avec le lot.

\pagerange{201}{211}
\section*{Endomorphismes de l'anneau différentiel simplicial (pages 201 à 211)}

On ne demande plus que $\varphi$ induise l'identité sur $\mathcal{A}_0 = S$.
$\mathcal{A}_*(S, S^+, t)$ est fonctoriel en $(S, S^+, t)$ pour les morphismes
$(u, \lambda)$, $u$ homomorphisme à puissances divisées et $\lambda \in S'$ avec
$u(t) = \lambda t'$, par $X_i \mapsto \lambda X_i$, $dX_i \mapsto \lambda dX_i$ ;
la composition est $(u', \lambda')(u, \lambda) = (u'u, u'(\lambda)\lambda')$, et
ces morphismes respectent la filtration naturelle $\mathrm{Fil}$\footnote{Dans
l'énoncé de la page 203 il écrit « $\lambda' \in S'$ », puis $u(t) = \lambda
t'$ et un but $\mathcal{A}_*(S, S'^+, t')$ ; nous lisons $\lambda \in S'$ et
$\mathcal{A}_*(S', S'^+, t')$.}. Si $t'$ est régulier, $\lambda$ est déterminé
par $u$ ; c'est le cas pour $S = k\{T\}$, $t = T$, lorsque $k$ est sans
torsion.

Soient $G_k$ le monoïde des $(u, \lambda)$ avec $u(T) = \lambda T$, $H_k$ celui
des endomorphismes $S_k$-linéaires, $\Omega_k$ celui de tous les
endomorphismes de $k$-algèbre différentielle simpliciale. On a
\[
G_k \cap H_k = \Bigl\{u_\lambda \Bigm| \lambda = 1 + \textstyle\sum_{i \geq 1}\lambda_iT^{(i)},\ (i+1)\lambda_i = 0\Bigr\} ,
\]
puisque $\lambda T = \sum_i (i+1)\lambda_iT^{(i+1)}$ ; c'est trivial si $k$ est
sans torsion, et si $k$ est une $\mathbb{Q}$-algèbre $\Omega_k$ est produit
semi-direct de $G_k \simeq G'_k = \operatorname{End}_{k\text{-alg.\ augm.}}(k\{T\})$ et du noyau.

Premier invariant : $\varphi_0 \in \operatorname{End}(S_k)$, donné par les
$\tau_i = \varphi_0(T^{(i)})$ avec $\tau_i\tau_j = c_{i,j}\tau_{i+j}$ ; chaque
$\varphi_n$ est $\varphi_0$-semi-linéaire. Second invariant : $\varphi_1$.
\textbf{Proposition} (page 209). \emph{Les homomorphismes d'anneaux
différentiels $S\{X_0\}[dX_0] \to S'\{X_0\}[dX_0]$ au-dessus de $\varphi_0$,
compatibles à l'augmentation $\alpha$, sont donnés par des $\xi_i = A_i +
B_i\,dX_0$ avec $A_i = A_1^{(i)} \in S'\{X_0\}^+$, via $\varphi_1(X_0^{(i)}) =
\xi_i$, $\varphi_1(dX_0) = d\xi_1$ ; graduations respectées si et seulement si
$B_i = 0$, puissances divisées si et seulement si $B_i = A_1^{(i-1)}B_1$.}
\emph{Remarque} : tout $\varphi_1$ s'écrit de façon unique $\psi_1 \circ
\overline{\varphi}_0$, où $\overline{\varphi}_0$ prolonge $\varphi_0$
canoniquement et $\psi_1$ est un $S'$-endomorphisme augmenté. La compatibilité
à $\beta$ donne $\beta(A_1) = \varphi_0(T)$ et $\varphi_0(T^{(i)}) =
\varphi_0(T)^{(i)}$ ; lorsque $\varphi_0(T) = T$, on retrouve $A_1 = X_0 + f$
avec $f = \sum_{i,j \geq 1} f_{ij}X_0^{(i)}X_1^{(j)}$\footnote{La page 211 écrit
$f = \sum X_0^{(i)}X_1^{(j)}$ sans coefficients, et sans rappeler la condition
$\varphi_0(T) = T$ sous laquelle la forme de $f$ est celle des pages 189 et
191.}. Le calcul de la page 201, qui développe $A_0$ en puissances de $X_0$,
appartient à ce même fil.

\pagerange{214}{216}
\section*{L'algèbre alternée libre à puissances divisées (pages 214 à 216)}

Sur une feuille marquée « Feuille à part ! » :

\textbf{Théorème.} \emph{Soient $(S, I)$ un anneau à idéal à puissances
divisées, $M^* = M^0 \oplus M^1$ un $S$-module gradué modulo $2$, $A(M^*) =
\Gamma^*_S(M^0) \otimes_S \Lambda^* M^1$, augmentée vers $S$, et $J =
\varepsilon^{-1}(I)$. a) Il existe une unique structure à puissances divisées,
au sens gradué, sur $J$ — c'est-à-dire sur ses éléments de degré pair, les
puissances divisées d'ordre $\geq 2$ d'un élément impair étant nulles —,
compatible à celle de $I$ et fonctorielle en $(S, M^*)$. b) Pour toute
$S$-algèbre alternée graduée modulo $2$, $B^*$, munie d'un idéal gradué à
puissances divisées $K^*$ compatible à $I$, la restriction}
\[
\operatorname{Hom}_{S\text{-alg.\ à p.d.}}\bigl((A(M^*), J), (B^*, K^*)\bigr) \xrightarrow{\ \sim\ } \operatorname{Hom}_{S\text{-mod.\ gr.}}(M^*, K^*)
\]
\emph{est bijective : $A(M^*)$ est l'algèbre alternée à puissances divisées
libre engendrée par $M^*$.}\footnote{Deux écarts. (i) La page formule a) pour
une structure à puissances divisées sur tout $J = J^0 + J^1$, prolongée aux
éléments impairs par $(x + u)^{(i)} = x^{(i)} + x^{(i-1)}u$. Elle découvre
elle-même, page 216, que l'additivité $(x+y)^{(i)} = \sum x^{(j)}y^{(i-j)}$ et la
règle $(\lambda x)^{(i)} = \lambda^i x^{(i)}$ ne valent alors que pour des
éléments qui « commutent strictement » ($x^1y^1 = 0$) : pour $u, v$ impairs,
$(u+v)^{(2)} = uv$ devrait être nul. La formulation correcte est celle de la
théorie de Cartan (Séminaire H.~Cartan, 1954--1955), où les puissances divisées
ne portent que sur les éléments de degré pair ; c'est celle que nous
énonçons. (ii) La vérification des axiomes est laissée inachevée page 216, à
l'axiome (4) $(x^{(i)})^{(j)} = \frac{(ij)!}{(i!)^j\,j!}x^{(ij)}$ ; elle est
reprise et achevée sur une autre feuille, pages 246 et 248. Pour la partie
paire $\Gamma_S(M^0)$ avec l'idéal $I$, c'est la théorie classique des
puissances divisées (Roby), que le dossier ne redémontre pas.}

\pagerange{218}{222}
\section*{Les quatre complexes (pages 218 à 222)}

Soient $A$ un anneau commutatif et $M$ un $A$-module ; on gradue par le degré
de $\Lambda^*M$.

(1) \emph{De Rham.} Sur $\operatorname{Sym}^*M \otimes \Lambda^*M$, l'unique
antidérivation $d$ nulle sur $\Lambda^1M$ et égale à l'isomorphisme canonique
sur $\operatorname{Sym}^1M$ ; $d^2 = 0$, $d : \operatorname{Sym}^p \otimes
\Lambda^q \to \operatorname{Sym}^{p-1} \otimes \Lambda^{q+1}$. C'est le complexe
de De Rham du fibré vectoriel $V(M)$. Si $M$ est plat et $A$ une
$\mathbb{Q}$-algèbre, sa cohomologie est $A$ en degré $0$ (lemme de Poincaré) ;
en caractéristique $p$, il y a « les opérations de Cartier »\footnote{Ce qu'on
appelle aujourd'hui l'isomorphisme de Cartier.}.

(2) \emph{Koszul.} L'unique antidérivation $\delta$ nulle sur
$\operatorname{Sym}^1M$, égale à l'isomorphisme $\Lambda^1M \to
\operatorname{Sym}^1M$ sur $\Lambda^1$ ; $\delta : \operatorname{Sym}^p \otimes
\Lambda^q \to \operatorname{Sym}^{p+1} \otimes \Lambda^{q-1}$. Si $M$ est plat,
sans condition de caractéristique, l'homologie est $A$ en degré $0$.

(3) \emph{De Rham à puissances divisées.} Sur $\Gamma^*M \otimes \Lambda^*M$,
l'unique antidérivation compatible aux puissances divisées, $d(x^{(i)}) =
x^{(i-1)}dx$, nulle sur $\Lambda^1M$ et égale à l'isomorphisme sur $\Gamma^1M$ :
c'est le complexe de De Rham à puissances divisées de $\Gamma^*M$. Si $M$ est
plat, sa cohomologie est $A$ en degré $0$, sans condition de caractéristique.

(4) \emph{Koszul à puissances divisées.} L'unique dérivation compatible aux
puissances divisées, nulle sur $\Gamma^1M$, induisant $\Lambda^1M \to \Gamma^1M$ ;
si $M$ est plat et $A$ de caractéristique $0$, l'homologie est $A$ en degré
$0$ ; en caractéristique $p$ « on doit avoir une histoire d'opérations de
Cartier ».

\emph{Relations} (page 220). Les quatre algèbres différentielles sont
fonctorielles en $(A, M)$, transforment sommes en produits tensoriels et
commutent aux limites inductives filtrantes. Le morphisme canonique
$\operatorname{Sym}^*M \otimes \Lambda^*M \to \Gamma^*M \otimes \Lambda^*M$ est
compatible à $d$ et à $\delta$, et c'est un isomorphisme en caractéristique
$0$. Le calcul de (3) se fait en se ramenant à $M$ libre de type fini (Lazard),
puis par Künneth à $M$ libre de rang $1$ ; celui de (4) se ramène à (2), qui se
traite comme (1). L'opérateur $\delta$ est un cas particulier du complexe de
Koszul d'une forme linéaire $\varphi : \mathcal{M} \to S$ : $H_0 =
S/\varphi(\mathcal{M})$, et pour $\mathcal{M}$ libre de type fini l'acyclicité en
degrés $> 0$ a lieu si une base donne une suite $S$-régulière, condition
nécessaire si $S$ est noethérien et $\varphi(\mathcal{M}) \subset
\operatorname{rad}S$.

\emph{L'opérateur d'Euler} (page 221). $\theta = d\delta + \delta d$, anticommutateur
de deux antidérivations, est une dérivation (compatible aux puissances divisées
dans le cas de $\Gamma$), égale à l'identité sur $\operatorname{Sym}^1$ (ou
$\Gamma^1$) et sur $\Lambda^1$ ; c'est donc la multiplication par le degré total
$n = p + q$ sur $\operatorname{Sym}^p \otimes \Lambda^q$ (ou $\Gamma^p \otimes
\Lambda^q$). Comme $d$ et $\delta$ conservent le degré total, chaque complexe se
scinde selon ce degré, et la composante de degré total $n$ est acyclique, pour
$d$ comme pour $\delta$, dès que $n \cdot 1_A$ est inversible — sans hypothèse
de platitude\footnote{La page dit « $H^n_{\mathrm{DR}}$ et $H^n_{\mathrm{Kosz}}$
sont nuls pour tout $n > 0$ », en mélangeant le degré extérieur et le degré
total ; c'est la composante de degré total $n$ que $\theta$ annule en
cohomologie.}. En particulier $\operatorname{Sym}^kM \xrightarrow{d}
\operatorname{Sym}^{k-1}M \otimes M$ est injectif si $k$ est inversible, car
$\delta d = k$ sur $\operatorname{Sym}^kM$.

\emph{Dualité.} Un accouplement $M \otimes M' \to A$ induit
$\operatorname{Sym}^*M \times \Gamma^*M' \to A$ et $\Lambda^*M \times \Lambda^*M'
\to A$, parfaits si $M$ est projectif de type fini et $M'$ son dual ; par ces
accouplements $d_M$ et $\delta_{M'}$, $\delta_M$ et $d_{M'}$, multiplication et
comultiplication sont transposés les uns des autres.

\pagerange{224}{225}
\section*{Deux feuilles volantes sur Dold--Puppe (pages 224 et 225)}

Deux feuillets d'une encre plus pâle, « paumés » dit la marge, appartiennent à
une rédaction numérotée qui n'est pas dans le dossier (ils renvoient à des
propositions 3, 4 et 5 absentes). La « Proposition 6 » exprime le foncteur
$\mathrm{ND}_\bullet$ des chaînes normalisées : $\mathrm{ND}_i(L_*) \simeq
\operatorname{Hom}_{k,*}(\Lambda^i\Phi_*, L_*)$, la différentielle provenant de
celle de $\Lambda^\bullet\Phi_*$, par la symétrie de Dold–Puppe et
$\mathrm{ND}(\Lambda^i\Phi_*) \simeq C^{i-1}_\bullet$ ; le corollaire donne la
version cosimpliciale $\mathrm{ND}^\bullet(L^*) \simeq \Lambda\Phi_* \otimes_{k,*}
L^*$. La page 225 en contient une première ébauche, biffée.

\pagerange{226}{232}
\section*{Calculs d'homotopie (pages 226 à 232)}

\textbf{Proposition 1.} \emph{$\Phi_*$ est homotope à $0$.} Par Dold–Puppe il
suffit que le complexe normalisé le soit : il est réduit à $k$ en degrés $0$ et
$1$, l'opérateur de bord étant l'identité.

\emph{Corollaire 1.1.} Pour tout foncteur $F$, $F(\Phi_*)$ est homotope à
$F(0)$ ; en particulier $\Lambda^i\Phi_*$, $\Gamma^i\Phi_*$,
$\operatorname{Sym}^i\Phi_*$ sont homotopes à $0$ pour $i \neq 0$.

\emph{Corollaire 1.2.} $\pi_i(\Psi_*) \simeq \pi_{i-1}(k_*)$ : $\Psi_*$ est un
$K(k,1)$, de complexe normalisé $k[1]$.

\emph{Corollaire 1.3.} $\Lambda^i\Psi_*$ est un $K(k, i)$. Dans le langage des
foncteurs dérivés de Dold–Puppe, c'est $\mathrm{L}\Lambda^i(k[1]) \simeq k[i]$ ;
on le voit aussi par la suite exacte $0 \to \Lambda^{i-1}\Psi_* \to
\Lambda^i\Phi_* \to \Lambda^i\Psi_* \to 0$, à terme médian contractile.

Les « mêmes arguments », ajoute la page, montreraient que $\Gamma^i\Psi_*$ et
$\operatorname{Sym}^i\Psi_*$ sont des $K(k,i)$, « nettement moins
économiques »\footnote{C'est faux en général. La formule de décalage de la page
361 donne $\mathrm{L}\operatorname{Sym}^i(k[1]) \simeq
\mathrm{L}\Lambda^i(k)[i]$, nul pour $i \geq 2$ puisque $\Lambda^i k = 0$ :
$\operatorname{Sym}^i\Psi_*$ est contractile pour $i \geq 2$. Et si $k$ est une
$\mathbb{Q}$-algèbre, $\Gamma^i \simeq \operatorname{Sym}^i$, donc
$\Gamma^i\Psi_*$ l'est aussi. L'argument par la suite exacte ne s'étend pas :
c'est pour $\Lambda$ que la suite $\Lambda^{i-1}\Psi \to \Lambda^i\Phi \to
\Lambda^i\Psi$ existe.}.

\emph{Corollaire 1.4.} Pour un $k$-module simplicial $A_*$, $\pi_\ell(A_*
\otimes \Lambda^i\Psi_*) \simeq \pi_{\ell-i}(A_*)$ : par Eilenberg–Zilber, le
complexe de $A_* \otimes B_*$ est équivalent à $A^!_* \otimes B^!_*$, et
$(\Lambda^i\Psi_*)^!$, complexe de projectifs d'homologie $k$ en degré $i$, est
homotopiquement équivalent à $k[i]$ ; c'est une « suspension additive ».

\emph{Corollaire 1.5.} Pour $C^{p,q}_* = \Gamma^p\Phi_* \otimes \Lambda^q
\Psi_*$, $\pi_\ell(C^{p,q}_*)$ vaut $k$ si $p = 0$ et $\ell = q$, et $0$ sinon.
En degré total $n$, les composantes $C^{n,0}_*, \ldots, C^{1,n-1}_*$ sont donc
flasques, et seule $C^{0,n}_*$ ne l'est pas : la ligne est le tronqué $\tau^{\leq
n}$ d'une résolution flasque de $k_*$, et
\[
C^{\bullet,\, \text{total } n}_{\mathrm{DRpd}}(X_*, k) \simeq \tau^{\leq n}\,\mathrm{R}\Gamma(X_*, k) , \qquad H^{p,q}_{\mathrm{DRpd}}(X_*, k) \simeq H^q(X_*, k) \quad (p \geq 0) ,
\]
ce qui redémontre le calcul de la page 8\footnote{La page indexe par (total
$p$, extérieur $i$) : $H^{p,i} \simeq H^i$ si $i \leq p$, $0$ si $i > p$.}.

\pagerange{232}{240}
\section*{La version complétée et filtrée (pages 232 à 240)}

Soient $S$ un anneau à idéal à puissances divisées $S^+$, $t \in S^+$, et
supposons $S$ séparé et complet pour la filtration par les $S^{[i]}$, adhérences
des idéaux engendrés par les $t^{[j]}$, $j \geq i$. On considère l'anneau
simplicial complété
\[
\widehat{A}^*_n = S\{\{X_0,\ldots,X_n\}\}/(\textstyle\sum X_i - t)_{\mathrm{pd}} \otimes_S \Lambda[dX_0,\ldots,dX_n]/(\textstyle\sum dX_i) ,
\]
de sorte que $\widehat{A}^*_* \simeq \widehat{A}^0_* \otimes_S \Lambda\Psi^S_*$
et $\pi_\ell(\widehat{A}^i_*) \simeq \pi_{\ell-i}(\widehat{A}^0_*)$. On filtre
$\widehat{A}^0_*$ par les idéaux fermés stables par puissances divisées
engendrés par la condition que les $X_i$ et $t$ soient de filtration $\geq 1$ :
$\mathrm{Fil}^p$ est formé des $\sum a_{i}X^{[i]}$ avec $a_i \in S^{[p - \sum
i_\alpha]}$. Avec $k = S/S^{[1]}$, et \emph{si $k\{T\} \to \mathrm{Gr}^*(S)$ est
un isomorphisme}, $\mathrm{Gr}^p(\widehat{A}^0_*) \simeq \Gamma^p_k(\Phi^k_*)$ est
contractile pour $p > 0$. Filtrant $\widehat{A}^*_*$ par $\mathrm{Fil}^p =
\sum_{i \leq p} \mathrm{Fil}^{p-i}(\widehat{A}^0_*) \otimes \Lambda^i\Psi_*$,
stable par la différentielle, on voit $\mathrm{Fil}^p(\widehat{A}^*_*)$ comme une
résolution de $\mathrm{Fil}^pS$ dont les termes de degré $< p$ sont flasques et
dont le terme de degré $p$ a pour seul groupe d'homotopie $\pi_p = k$. D'où,
pour $q \leq p$, un carré commutatif
\[
\begin{array}{ccc}
H^i\bigl(\mathrm{Fil}^p C^*_{\mathrm{DR}}(\mathcal{X}; S, t, \gamma)\bigr) & \xrightarrow{\ \sim\ } & H^i(\mathcal{X}, \mathrm{Fil}^pS) \\
\downarrow & & \downarrow \\
H^i\bigl(\mathrm{Fil}^q C^*_{\mathrm{DR}}(\mathcal{X}; S, t, \gamma)\bigr) & \longrightarrow & H^i(\mathcal{X}, \mathrm{Fil}^qS)
\end{array}
\]
dont la flèche du haut est un isomorphisme pour $i \leq p$, celle du bas pour
$i \leq q$.

Sans l'hypothèse sur $\mathrm{Gr}^*(S)$, « on ne peut rien dire ». Pour $t =
0$, $S \simeq k$ et $\mathrm{Gr}^p \simeq \Gamma^p_k(\Psi_*)$, qui n'est flasque
pour aucun $p$ si $k \neq 0$. Pour $(\mathbb{Z}_p, p)$, la filtration
à puissances divisées définie par $p$ est cofinale à la filtration $p$-adique si
$p \neq 2$ ; pour $p = 2$ elle est constante, $\mathrm{Fil}^i\mathbb{Z}_2 =
2\mathbb{Z}_2$ pour $i \geq 1$, et $\mathrm{Gr}^*(S) \simeq \mathbb{F}_2$ : cas
dégénéré. Le calcul de la page 240 le justifie : avec la formule de
Legendre, $v_p(p^n/n!) = \frac{(p-2)n + s_p(n)}{p-1}$, et
$\gamma_{p,i} = \inf_{n \geq i} v_p(p^n/n!)$ vaut $1$ si $p = 2$ (les puissances
de $2$ ont un seul chiffre) et est $> \frac{p-2}{p-1}\,i \to +\infty$ si $p
\neq 2$. Si $i$ se termine par $s$ chiffres $p - 1$, on a $v_p(p^{(i+1)}) -
v_p(p^{(i)}) = 1 - s$, donc
\[
v_p\bigl(p^{(i)}\bigr) \leq v_p\bigl(p^{(i+1)}\bigr) \iff s \leq 1 \iff p^2 \nmid i+1 ,
\]
et $\gamma_{p,i} = \gamma_{p,i+1}$ dès que $p \mid i+1$\footnote{La page
conclut « $\iff s = 0 \iff i+1 \not\equiv 0 \pmod p$ » : en comparant les sommes
de chiffres de $i$ et de $i+1$, elle omet l'unité ajoutée au chiffre $i_s$. Pour
$p = 3$, $i = 2$, on a $v_3(3^2/2!) = 2 = v_3(3^3/3!)$, avec $3 \mid i+1$. La
seconde affirmation, « il suffit que $i + 1 \equiv 0 \pmod p$ » pour
$\gamma_{p,i} = \gamma_{p,i+1}$, est exacte.}.

\pagerange{243}{260}
\section*{Les catégories $\mathcal{M}_N$ et les modules négligeables (pages 243 à 260)}

\subsection*{Adjoints de l'inclusion}

Soit $S$ un anneau commutatif $\mathbb{Z}$-gradué à degrés $\geq 0$, $S^0 = k$,
$\mathcal{M}$ la catégorie des $S$-modules gradués, $\mathcal{M}_N$ celle des
modules nuls en degrés $< N$ ($\mathcal{M}_N = \mathcal{M}_0[-N]$),
$\mathcal{N}$ celle des $k$-modules, et $i : \mathcal{N} \to \mathcal{M}_0$, $M
\mapsto M \otimes_k S$. L'inclusion $i_{N,N'} : \mathcal{M}_{N'} \hookrightarrow
\mathcal{M}_N$ ($N' \geq N$) commute aux limites ; elle a un adjoint à droite,
la troncature $\tau_{N'}$ (sous-module $\bigoplus_{i \geq N'} M_i$), et un
adjoint à gauche $q_{N'}(M) = M/\langle M_i,\ i < N'\rangle$. Ce sont des
« foncteurs de localisation » ; $q$ est exact à droite, $\tau$ exact et commute
à toutes les limites, donc est le quotient par la sous-catégorie épaisse des
modules nuls en degrés $\geq N'$, et a encore un adjoint à droite :
\[
\rho_{N'}(P)_i = \operatorname{Hom}^i_{\mathcal{M}}\bigl(\tau_{N'-i}S, P\bigr) , \qquad \rho_{N,N'} = \tau_N\,\rho_{N'} .
\]
On a donc la chaîne d'adjoints $q_{N'} \dashv i_{N'} \dashv \tau_{N'} \dashv
\rho_{N'}$, $i$ et $\rho$ pleinement fidèles\footnote{Le foncteur $\rho$ de
ces pages est le $\varphi$ des pages 26, 256 et 290 ; nous gardons $\varphi$
ailleurs.}.

Le foncteur $\tau_N \circ i : \mathcal{N} \to \mathcal{M}_N$ a pour adjoint à
droite $k \circ \rho_N$ ($k$ = composante de degré $0$), et sont équivalents
(page 251) : a) $\tau_N \circ i$ est pleinement fidèle ; b) $P
\xrightarrow{\sim} (\rho_N\tau_N(P \otimes_k S))_0$ ; c) $k \circ \rho_N$ est
une localisation ; et aussi c') pour $0 \leq \alpha \leq N-1$, $P \otimes
S_\alpha \xrightarrow{\sim} \operatorname{Hom}^\alpha(\tau_{N-\alpha}S, P
\otimes_k S)$, soit $\operatorname{Hom}^\alpha(S/\tau_{N-\alpha}S, P \otimes_k S)
= 0$ — c'est le théorème des pages 26 à 31.

\subsection*{Algèbres différentielles graduées à puissances divisées}

Les feuilles 246 et 248, écrites tête-bêche, fixent les axiomes d'une algèbre
$(A^*, \cdot, d, J, \gamma)$ : groupe gradué à degrés $\geq 0$, algèbre
associative unitaire alternée ($xy = (-1)^{ij}yx$, $x^2 = 0$ en degré impair),
différentielle de degré $+1$ dérivation graduée, idéal $J = J^0 + A^+$,
puissances divisées sur $J$ avec $x^{(\alpha)} \in A^{\alpha i}$, $x^{(\alpha)} =
0$ pour $x$ impair et $\alpha \geq 2$, et $d\,x^{(\alpha)} = x^{(\alpha-1)}dx$.
Elles vérifient, pour $z = x + y$ ($x \in J^0$, $y$ impair), $z^{(\alpha)} =
x^{(\alpha)} + x^{(\alpha-1)}y$ : l'additivité « OK si $yy' = 0$ », la règle
$(\lambda z)^{(\alpha)} = \lambda^\alpha z^{(\alpha)}$ « OK si $\eta y = 0$ »,
et, sans condition, $z^{(p)}z^{(q)} = c_{p,q}z^{(p+q)}$ et $(z^{(\alpha)})^{(\beta)}
= \frac{(\alpha\beta)!}{\beta!(\alpha!)^\beta}z^{(\alpha\beta)}$. C'est
l'achèvement de la vérification laissée en suspens page 216.

La page 250 écrit le complexe $C^*(X, (A, \pi, \gamma))$ : pour un anneau $A$ à
puissances divisées et $\pi$ dans l'idéal, $A\{(X_i, dX_i)\}/(\sum X_i - \pi,
\sum dX_i)$, contravariant en $X$ et covariant en $(A,\pi,\gamma)$ pour les
couples $(\varphi, \lambda)$ avec $\varphi(\pi) = \lambda\pi'$, $X_i \mapsto
\lambda X_i$. La page 252, tête-bêche elle aussi, est une suite de formules
sans texte autour des $\Gamma^*_{(A,\pi)}(J^I)$, de la flèche $H^{2i}(X, J) \to
H^{2i\alpha}(X, \Gamma^\alpha J)$ et de $i!\,x^{(i)} = x^i$ ; elle ne se lit pas
comme un argument.

\subsection*{Modules négligeables}

Un module est \emph{$N$-négligeable} si $M_i = 0$ pour $i \geq N$,
\emph{$\infty$-négligeable} s'il l'est pour un $N$, \emph{négligeable} si
chacun de ses éléments est annulé par un $\tau_NS$, c'est-à-dire s'il est
limite inductive de modules $\infty$-négligeables ; pour un module de type
fini, négligeable équivaut à $\infty$-négligeable. Ces classes forment des
sous-catégories épaisses, avec $\mathcal{M}/\mathrm{Nég}_N \simeq
\mathcal{M}_N$. Notons $\mathcal{M}_\omega = \mathcal{M}/\mathrm{Nég}$ et
$\tau_\omega$, $\varphi_\omega$ le foncteur quotient et son adjoint ; la page
253 note que $\varinjlim \mathcal{M}_N \simeq \mathcal{M}/\mathrm{Nég}_\infty$
n'est pas une catégorie à limites inductives filtrantes exactes\footnote{C'est
le cadre de ce qu'on appelle aujourd'hui une catégorie quotient de Serre, ou
localisation au sens de Gabriel.}.

\emph{Proposition} (page 256) : $\varinjlim_N \varphi_N\tau_N P
\xrightarrow{\sim} \varphi_\omega\tau_\omega P$. On teste sur les modules de
présentation finie $M$ : les deux membres sont $\varinjlim
\operatorname{Hom}(\tau_NM, P)$ et $\varinjlim_{M' \subset M,\ M/M' \text{
négl.}} \operatorname{Hom}(M', P)$, et pour $M$ de type fini les deux systèmes
de sous-objets sont cofinaux. En particulier $(\varphi_\omega\tau_\omega P)_i
\simeq \varinjlim_N \operatorname{Hom}^i(\tau_NS, P)$, et $P$ est dans l'image
essentielle de $\varphi_\omega$ si et seulement si a) $P$ n'a pas d'élément
négligeable non nul, et b) tout homomorphisme de degré $i$ d'un $\tau_NS$ dans
$P$ se prolonge à $S$ — toute extension de $(S/\tau_NS)[-i]$ par $P$ est triviale.

\emph{Foncteurs dérivés} (page 260). De $\operatorname{Hom}_{\mathcal{M}}(M,
\varphi_N\tau_NP) = \operatorname{Hom}(i_N\tau_NM, P)$ et de l'exactitude de
$\tau_N$ vient, pour $M$ projectif,
\[
\operatorname{Hom}\bigl(M, \mathrm{R}\varphi_N(\tau_NP)\bigr) \simeq \mathrm{R}\operatorname{Hom}(\tau_NM, P) , \qquad \operatorname{Hom}\bigl(M, \mathrm{R}^n\varphi_N(\tau_NP)\bigr) \simeq \operatorname{Ext}^n(\tau_NM, P) .
\]

\pagerange{262}{276}
\section*{Reconstruire un complexe à partir de son ombre (pages 262 à 276)}

Avec $M = S[-i]$ : $\bigl(\mathrm{R}^n\varphi_N(\tau_NP)\bigr)_i \simeq
\operatorname{Ext}^n\bigl((\tau_{N-i}S)[-i], P\bigr)$, et de même pour
$\varphi_\omega$ en passant à la limite sur $N$.

Soient $C$ une catégorie abélienne et $i : C \to \mathcal{M}_{N_0}$ un foncteur
exact tel que $i \simeq \varphi\tau\, i$. La page en conclut
$(\mathrm{R}^n\varphi)(\tau\, i(M)) = 0$ pour $n > 0$, en écrivant $\mathrm{R}^*i
\simeq \mathrm{R}^*(\varphi\tau i) \simeq (\mathrm{R}^*\varphi)(\tau i)$ « car
$\tau i$ est exact ». Appliqué à $C = \mathrm{Mod}(k)$, $i(M) = M \otimes_k S$ —
qui est bien dans l'image essentielle, par le théorème de la page 26 —, cela
donnerait l'énoncé encadré
\[
\operatorname{Ext}^n_{\mathcal{M}}\bigl((\tau_{N-i}S)[-i], M \otimes_k S\bigr) = 0 \qquad (n > 0,\ i \geq 0) ,
\]
le corollaire annoncé page 26. La justification ne suffit pas telle
quelle\footnote{Pour composer les foncteurs dérivés, $\mathrm{R}(\varphi \circ
G) \simeq \mathrm{R}\varphi \circ G$, il faut que $G = \tau \circ i$ envoie les
injectifs de $C$ sur des objets $\varphi$-acycliques, et l'exactitude de $G$ ne
le donne pas : $\tau$ préserve les injectifs (son adjoint à gauche est exact),
mais $i$, adjoint à \emph{gauche} du foncteur exact « composante de degré $0$ »,
préserve les projectifs et non les injectifs. Nous n'avons pas établi
l'annulation par un autre chemin, et la laissons comme un énoncé du dossier,
non démontré. Ce qui suit en dépend.}.

Sous cette annulation, la reconstruction suit (pages 266 à 272). Soient $K^\bullet$
un complexe de $C$, $L^\bullet$ un complexe de $\mathcal{M}$ et $\alpha :
i(K^\bullet) \to L^\bullet$ dans $D(\mathcal{M}_{N_0})$ tel que
$\tau_\omega\alpha$ soit un isomorphisme dans $D(\mathcal{M}_\omega)$ — noyau et
conoyau des $H^n(\alpha)$ négligeables. Alors $i(K) \simeq
\mathrm{R}\varphi_\omega(\tau_\omega L^\bullet)$, et si $i$ est pleinement fidèle
avec un adjoint à droite $k$ exact, $k \circ i \simeq \mathrm{id}$, on obtient
dans $D(C)$
\[
K^\bullet \simeq \mathrm{R}(k \circ \varphi_\omega)\bigl(\tau_\omega L^\bullet\bigr) :
\]
\emph{$K^\bullet$ se reconstitue à partir de son ombre $\tau_\omega(i(K^\bullet))
\simeq \tau_\omega(L^\bullet)$}. En termes de $\sigma = \tau_\omega \circ i : C \to
\mathcal{M}_\omega$, pleinement fidèle et exact, d'adjoint à droite $\psi =
k\varphi_\omega$ (« peut-être pas exact »), avec $\psi\sigma \simeq
\mathrm{id}$, la page 272 énonce : $D^+(C) \to D^+(\mathcal{M}_\omega)$ est
pleinement fidèle, compatible à la multiplication\footnote{Une tentative de
simplification, $K \simeq k\,\varphi_\omega(\tau_\omega L)$ sans dériver, est
encadrée et barrée page 268 ; un calcul par la suite spectrale $E_2^{pq} =
\mathrm{R}^q\psi(H^p(L))$, page 270, est aussi traversé de traits.}. C'est le
« en raffinant un peu » de la page 10 : la catégorie dérivée des $k$-modules se
plonge dans celle des ombres.

Les pages 274 et 276 cherchent à écrire explicitement $i \to
\varphi_{0,N}\tau_{N,0}\,i$ comme un noyau, $M \hookrightarrow
M^{[N,+\infty[} \rightrightarrows M^{[N,+\infty[ \times \mathbb{N}}$, avec les
relations $c_{\ell,k}x_k - c_{\ell,k-i}x_{k+\ell}$, et définissent, dans une
catégorie abélienne $k$-linéaire, les objets $\underline{\operatorname{Hom}}(M,
X)$ et $M \otimes_k X$ par leurs propriétés universelles ; la page 276 est un
brouillon sur l'exactitude d'un foncteur, sans conclusion.

\pagerange{278}{284}
\section*{Le produit cotensoriel (pages 278 à 284)}

La page 278 reprend le point de vue des coalgèbres : $S = k[U]$ ($U^n$ de degré
$-n$), d'algèbre duale $k\{\{T\}\}$, le foncteur $i_* : M \mapsto M \otimes_k S$,
et la cotroncature $\sigma_N(P) = P/\tau_{N+1}(P)$ ; un théorème de pleine
fidélité est énoncé et le calcul s'arrête sur
« $\operatorname{Hom}_{\mathcal{M}}(S, \tau^NN \otimes_kS) \simeq N$ ?? ». Les pages
280 à 284 posent, pour deux comodules $C$, $C'$ sur la coalgèbre $S$,
\[
P = C \,\square_S\, C' = \operatorname{Ker}\bigl(C \otimes C' \rightrightarrows C \otimes S \otimes C'\bigr) ,
\]
le produit tensoriel « au-dessus de $S$ » des comodules, dual de
$V(C) \boxtimes_{V(S)} V(C')$\footnote{La page écrit le terme de droite $C
\otimes C' \otimes S$, puis biffe le $\otimes S$ final ; c'est ce qu'on appelle
aujourd'hui le \emph{produit cotensoriel}. Le passage aux intersections
$P_\infty = \bigcap P_i$ d'une suite décroissante, et la question de savoir si
$P \to P \otimes S$ est bien le coproduit, restent des calculs sans
conclusion. La moitié inférieure de la page 280, tête-bêche et raturée, porte
sur les coalgèbres à puissances divisées et se lit mal.} ; la question
encadrée est $(M \otimes_k S) \boxtimes_S C \simeq M \otimes_k C$ ?, et la page 284
écrit qu'un accouplement $V(C) \times V(C') \to V(C'')$ bilinéaire par rapport
à $V(S)$ revient à un morphisme de coalgèbres $C'' \to C \,\square_S\, C'$.

\pagerange{286}{298}
\section*{Modules gradués : la mise au propre (pages 286 à 298)}

$S$ commutatif $\mathbb{Z}$-gradué à degrés $\geq 0$, $k = S_0$,
$\mathcal{M}$ les $S$-modules gradués, $M[i]$ la translation. Pour tout $N$, on
pose $\mathcal{M}_N$ (modules nuls en degrés $< N$) et $\mathcal{M}^N$ (nuls en
degrés $> N$) ; $\mathcal{M}_N = \mathcal{M}_0[-N]$, $\mathcal{M}^N =
\mathcal{M}^0[-N]$, et toute la donnée équivaut au couple $(\mathcal{M}_0,
\mathcal{M}^0)$.

\textbf{Proposition 1.} \emph{L'inclusion $i_N : \mathcal{M}_N \hookrightarrow
\mathcal{M}$ s'insère dans une suite de quatre foncteurs adjoints $q_N \dashv
i_N \dashv \tau_N \dashv \varphi_N$, avec $\tau_N(P) = \bigoplus_{i \geq N}P_i
\subset P$ et}
\[
(\varphi_N Q)_i = \operatorname{Hom}^i_{\mathcal{M}}\bigl(\tau_{N-i}(S), Q\bigr) = \operatorname{Hom}^*_{\mathcal{M}}\bigl(\tau_{N-i}(S)[-i], Q\bigr)_0 ,
\]
\emph{la structure de module venant de la multiplication $\tau_{N-i-k}S \to
\tau_{N-i}S$ par $u \in S_k$, transposée.} La coünité $\tau_N\varphi_N Q
\xrightarrow{\sim} Q$ est un isomorphisme ; l'unité $P \to \varphi_N\tau_NP$
vient de $P_i \simeq \operatorname{Hom}^i(S, P) \to \operatorname{Hom}^i(\tau_{N-i}S,
\tau_NP)$, tout homomorphisme de degré $i$ de $\tau_{N-i}S$ dans $P$ tombant dans
$\tau_NP$. Enfin $q_N(P) = P/\langle P_i,\ i < N\rangle$, « qu'on n'utilisera
pas ».

\emph{Corollaire.} $i_N$ et $\tau_N$ commutent à toutes les limites, $q_N$ aux
limites inductives, $\varphi_N$ aux limites projectives ; $i_N$ et $\varphi_N$
sont pleinement fidèles, $q_N$ et $\tau_N$ des localisations ; $\tau_N$ est la
localisation par la sous-catégorie épaisse $\mathcal{M}^{N-1}$.

La page 294 généralise à $\mathcal{M}(C,S)$, les $S$-modules gradués dans une
catégorie abélienne $C$ : l'adjoint $\varphi_N$ existe si $C$ ou $C^\circ$ est
une catégorie de Grothendieck, le foncteur $P \mapsto
\operatorname{Hom}(\tau_NP, Q)$ transformant limites inductives en limites
projectives, et la même formule le décrit\footnote{L'étape intermédiaire porte
« justifier » au-dessus du signe d'égalité.} ; par dualité, $\mathcal{M}^N(C,S)^\circ
\simeq \mathcal{M}_{-N}(C^\circ, S)$, et $i^N$ s'insère dans la suite
$\varphi^N \dashv \tau^N \dashv i^N \dashv q^N$. Le feuillet 298, sans prose,
reprend les homomorphismes homogènes entre tronqués des pages 36 à 40 et la
relation $c_{i,\ell}\,\xi_{i+\ell} = c_{i-d,\ell}\,\xi_i$ pour les homomorphismes
de degré négatif.

\pagerange{300}{309}
\section*{Notations simpliciales. Constructions universelles (pages 300 à 309)}

Une chemise de sa main (page 300) ouvre une série paginée de 1 à 17.
$\Delta$ est la catégorie des ordinaux $\Delta^n = [0,n]$, $\widehat{\Delta} =
\mathcal{S}s = \mathrm{Ens}_*$ ; pour une catégorie $\mathcal{A}$, $\mathcal{A}_*
= \underline{\operatorname{Hom}}(\Delta^\circ, \mathcal{A})$ (objets simpliciaux),
$\mathcal{A}^* = \underline{\operatorname{Hom}}(\Delta, \mathcal{A})$ (objets
cosimpliciaux), et $(\mathcal{A}^\circ)^* = (\mathcal{A}_*)^\circ$\footnote{Il
écrit $\underline{\operatorname{Hom}}(\Delta^\circ, \mathrm{Ens})$ et
$\underline{\operatorname{Hom}}(\Delta, \mathrm{Ens})$ là où il faut
$\mathcal{A}$. La marge précise : la notation $\widehat{\Delta}$ quand on veut
garder à l'esprit la nature algébrique très particulière de ce topos,
$\mathcal{S}s$ quand on veut l'oublier au profit de sa signification
topologique.}. Le programme est énoncé d'emblée : étudier les objets
simpliciaux — « l'équivalent combinatoire des espaces topologiques » — et leurs
types d'homotopie via des invariants covariants ou contravariants, « dans
l'esprit de l'algèbre universelle », avec leurs structures et les opérations
qu'on peut définir sur eux.

Les composantes s'écrivent $K_{[n]}$ et non $K_n$, réservé aux complexes de
chaînes ; une dualité $v : \mathcal{A}^\circ \simeq \mathcal{B}$ (typiquement les
modules projectifs de type fini sur $k$ commutatif) fait « monter et
descendre les indices », $K^{[n]} = K_{[n]}^\vee$. Le foncteur pleinement fidèle
$\Delta \hookrightarrow \widehat{\Delta}$, $n \mapsto \Delta^n_*$, est l'\emph{objet
cosimplicial universel} $\Delta^*_*$, avec $\operatorname{Hom}(\Delta^n_*, X_*)
\simeq X_{[n]}$. Les objets à plusieurs indices s'écrivent $K^*{}_*$, $K^{**}$,
$K_{**}$, avec un symbole ${}^\sigma K$ pour le symétrisé, afin de ne pas
brouiller la convention de montée des indices.

\textbf{Proposition 1.} \emph{Si $\mathcal{A}$ a des limites inductives,
$F \mapsto F(\Delta^*)$ est une équivalence de la catégorie des foncteurs
$\mathcal{S}s \to \mathcal{A}$ commutant aux limites inductives sur
$\mathcal{A}^*$ ; dualement pour $\mathcal{B}$ à limites projectives et les
foncteurs $\mathcal{S}s^\circ \to \mathcal{B}$ transformant limites inductives
en limites projectives, avec $\mathcal{B}_*$.} L'inverse est $X_* \mapsto
\varinjlim_{\Delta^n \in \Delta/X_*} K^{[n]}$ ; c'est vrai pour toute petite
catégorie, et « bien connu »\footnote{C'est l'extension de Kan à gauche le long
du plongement de Yoneda, ou le couple nerf–réalisation. La page écrit $K_*$,
$K_{[n]}$ pour l'objet cosimplicial, et $\varphi$ pour $\psi$ dans la version
duale.}. Le yoga : les invariants covariants commutant aux limites inductives
sont des objets cosimpliciaux de $\mathcal{A}$ ; les invariants contravariants
correspondants sont des objets simpliciaux de $\mathcal{B}$, et pour
$\mathcal{B} = \mathrm{Ens}$ des ensembles simpliciaux $K_*$, représentant $X_*
\mapsto \operatorname{Hom}(X_*, K_*)$.

Puis la réflexion de méthode qui oriente tout le dossier. Ces invariants obtenus
« à coups de foncteurs représentables » sont de nature trop « grosse » et
« fruste » pour être intéressants directement — ce ne sont pas des invariants du
type d'homotopie. On devra en extraire d'autres, plus subtils, « qui ne seront
pas de nature si simple ». « Notre propos sera de mettre sur les invariants gros
[\ldots] une machinerie de structures supplémentaires raisonnables, et dans les
catégories de structures de ce type (complexes de chaînes ou de cochaînes,
algèbres cosimpliciales etc.) d'inverser les flèches qui sont obtenues à partir
d'équivalences d'homotopie [\ldots] L'invariant ainsi obtenu, plus grossier bien
sûr, mais plus subtil dans sa définition, et moins redondant, exprime donc de
façon plus ou moins complète le type d'homotopie, ou en saisit de façon
adéquate tel ou tel aspect. »

Pour $\mathcal{A}$ additive karoubienne, la théorie de Dold–Puppe donne des
équivalences $\mathrm{ND} : \mathcal{A}_* \rightleftarrows \mathcal{A}_\bullet
: \mathrm{DP}$ avec les complexes de chaînes, et dualement avec les complexes de
cochaînes ; la même lettre désigne alors quatre avatars $K^*$, $K_*$,
$K^\bullet$, $K_\bullet$\footnote{C'est la correspondance de Dold–Kan ; $\mathrm{ND}$
est le complexe normalisé (« non dégénéré »), $\mathrm{DP}$ son inverse.}, et les
objets mixtes sont des bicomplexes. Remarque essentielle : les bicomplexes sont
à différentielles qui \emph{commutent}, contrairement aux conventions de
Cartan–Eilenberg, avec un signe dans la différentielle totale ; il « n'a pas eu
encore besoin vraiment » d'une convention de degré pour les bicomplexes mixtes.

\pagerange{310}{317}
\section*{L'invariant additif universel (pages 310 à 317)}

\textbf{Proposition 2.} \emph{Si $\mathcal{A}$ est additive à limites
inductives, les foncteurs $\mathcal{S}s \to \mathcal{A}$ commutant aux limites
inductives équivalent aux complexes de cochaînes $\mathcal{A}^\bullet$ ; si
$\mathcal{B}$ est additive à limites projectives, les foncteurs $\mathcal{S}s^\circ
\to \mathcal{B}$ qui transforment limites inductives en limites projectives
équivalent aux complexes de chaînes $\mathcal{B}_\bullet$.}\footnote{L'indice
de la seconde ligne se lit « comm.\ aux $\varinjlim$ », avec doute ; c'est la
transformation des limites inductives en limites projectives qui convient.}

Pour $\mathcal{B} = \mathrm{Ab}$ : si $K_\bullet$ correspond à $F$, et
$\mathbb{C}_*(X_*) = \mathbb{Z}(X_*) = (n \mapsto \mathbb{Z}^{(X_{[n]})})$ est
l'abélianisé de $X_*$, $\mathbb{C}_\bullet(X_*)$ son complexe normalisé,
\[
F(X_*) \simeq \operatorname{Hom}_{\mathrm{Ab}_*}\bigl(\mathbb{C}_*(X_*), K_*\bigr) \simeq \operatorname{Hom}_{\mathrm{Ab}_\bullet}\bigl(\mathbb{C}_\bullet(X_*), K_\bullet\bigr) .
\]
Deux conclusions. 1° Ces invariants ne dépendent que de $\mathbb{C}_\bullet(X_*)$,
« l'invariant additif universel » ; vus comme foncteurs de $C_\bullet$, ce sont
exactement les représentables. Une structure plus riche sur $F$ — une
multiplication — n'est pas reconstruite par $C_\bullet$ seul : il faut tenir
compte de structures supplémentaires de $C_\bullet$ « qui peuvent s'expliciter en
se rappelant que $C_\bullet$ provient d'un $X_*$ », et « l'étude systématique de
ces structures » est, dit-il, « le contenu essentiel de nos réflexions ». Une
structure de $k$-module, en revanche, se lit sur $C_\bullet \otimes k$. 2° Il y a
une solution universelle, $F_0(X_*) = \mathbb{C}_\bullet(X_*)$ à valeurs dans
$\mathcal{B}_0 = (\mathrm{Ab}_\bullet)^\circ$.

\textbf{Proposition 3} (« nouvelle forme de l'énoncé du théorème de
Dold–Puppe »). \emph{Pour toute catégorie additive $\mathcal{B}$ à limites
projectives, $\varphi \mapsto \varphi \circ F_0$ est une équivalence entre les
foncteurs $\mathcal{B}_0 \to \mathcal{B}$ et les foncteurs $\mathcal{S}s^\circ
\to \mathcal{B}$ qui commutent aux limites au sens ci-dessus.} Sous la forme
équivalente avec $\widehat{\Delta}_{\mathrm{ab}}$ et l'abélianisation
$G_0 : X \mapsto \mathbb{Z}(X)$, l'énoncé vaut pour toute petite catégorie à la
place de $\Delta$ : un foncteur $\widehat{\Delta}_{\mathrm{ab}} \to \mathcal{A}$
commutant aux limites inductives a un adjoint à droite, et se lit sur
$\underline{\operatorname{Hom}}(\Delta, \mathcal{A})$\footnote{C'est dire que la
catégorie des préfaisceaux abéliens sur $\Delta$ est la complétion libre de
$\Delta$ par limites inductives dans le monde additif, et que l'abélianisation
est l'objet universel. La démonstration de la page 313 est esquissée (« on a
pratiquement terminé »).}. En variante : l'objet cosimplicial $n \mapsto
\mathbb{Z}(\Delta^n_*)$ de $\mathrm{Ab}_*$ est universel pour les objets
cosimpliciaux des catégories additives à limites inductives.

Du point de vue de l'algèbre universelle : les invariants additifs s'expriment
par des objets simpliciaux ou des complexes, et les « opérations sur les
invariants » $\Omega$ par des objets $L_*$ ou $L_\bullet$ de $\mathrm{Ab}_*$ ou
$\mathrm{Ab}_\bullet$, avec $\Omega(F) \simeq \operatorname{Hom}(L_*, K_*)$, et
dualement $G(X_*) \simeq \mathbb{C}_*(X_*) \otimes K^*$ pour les covariants.
Remarques. a) Les mêmes $L$ servent, contravariants dans un cas, covariants dans
l'autre. b) La question des constructions possibles est ainsi résolue « de façon
quasi-tautologique », parce qu'on a admis des catégories de valeurs très
générales ; si l'on se restreint à un type $L$ de limites, on doit trouver la
sous-catégorie pleine $U_L$ engendrée par les $\mathbb{C}_*(\Delta^n)$ par ces
limites — « il doit résulter des principes généraux (que je n'ai pas
détaillés) » ; on le vérifie pour $L = \varnothing$, sommes finies. c) Pour
$\mathcal{B} = \mathrm{Ab}$, $F(X_*) \simeq \Omega_*(\mathbb{C}_*(X_*)^\circ)$.
Tout cela vaut pour une petite catégorie quelconque, et dans le contexte non
additif\footnote{La liste passe de c) à e) : il n'y a pas de d) sur ces pages.
Les fins des pages 311 et 316 sont très surchargées, avec des insertions
superposées dont l'ordre est incertain ; la lecture de b) et c) est mince.}.

e) \emph{Le rôle de Dold–Puppe.} L'avantage des complexes sur les objets
simpliciaux est que la description d'un objet y est « nettement plus simple — on
se perd facilement dans la description des opérations simpliciales ! » ; $L_*$
apparaît comme « une sorte de version pléthorique de $L_\bullet$ ». « Néanmoins,
dès que $\mathrm{Ab}_{k*}$ et $\mathrm{Ab}_{k\bullet}$ ont chacun leur structure
multiplicative — qui ne se correspondent \emph{pas} par DP et ND —, celle de
$\mathrm{Ab}_{k*}$ est plus simple à beaucoup d'égards, et s'introduit d'ailleurs
de bien des façons par la suite de façon absolument impérieuse », et les
opérations $\Lambda^i$, $\otimes^i$, $\operatorname{Sym}^i$ « manquent purement et
simplement » sur les complexes, sauf à les y transporter par DP\footnote{C'est
le fait, aujourd'hui standard, que la correspondance de Dold–Kan n'est pas
monoïdale : elle ne l'est qu'à homotopie près, par les applications
d'Eilenberg–Zilber et d'Alexander–Whitney (pages 356 à 359).}.

\pagerange{319}{324}
\section*{Gribouillis : coalgèbres et puissances divisées (pages 319 à 324)}

Une chemise bleue (page 318) porte « Gribouillis vrac » et « Au fil des
jours… ». Pour une coalgèbre $A$ et $F_A(X) = \operatorname{Hom}(A, X)$, il
cherche une transformation naturelle $\Gamma^*F_A(X) \to F_A(\Gamma^*X)$ et
montre que les données suivantes se déterminent (pages 319 à 322, puis 328) :
(1) une quasi-commutation de $F_A$ aux $\Gamma^p$, compatible à la structure
multiplicative $\beta_{pq}$ de $\Gamma$ et à la composition $\gamma_{pq} :
\Gamma^{pq} \to \Gamma^p\Gamma^q$ ; (2) des $u_p \in \operatorname{Hom}(A,
\Gamma^pA)$ tels que
\[
u_0 = 1,\quad u_1 = \mathrm{id}, \qquad u_p\,u_q = c_{p,q}\,u_{p+q} , \qquad \Gamma^p(u_q) \circ u_p = \gamma^A_{pq} \circ u_{pq} ,
\]
c'est-à-dire un homomorphisme d'algèbres à puissances divisées $k\{T\} \to
\operatorname{Hom}(A, \Gamma^*A)$ avec $\Gamma^p(u(T^{(q)}))\,u(T^{(p)}) =
\gamma^A_{pq}\,u(T^{(pq)})$ ; (3) des applications non linéaires $f \mapsto
f^{(p)}$ de $\operatorname{Hom}_k(A, M)$ dans $\operatorname{Hom}_k(A,
\Gamma^pM)$, fonctorielles en $M$, avec $(f+g)^{(p)} = \sum f^{(i)}g^{(j)}$ ;
(4)–(6) des puissances divisées fonctorielles sur $\operatorname{Hom}_k(A, J)$
pour tout anneau à idéal à puissances divisées $(B, J)$. L'additivité se vérifie
sur le cas universel $f = \mathrm{pr}_1$, $g = \mathrm{pr}_2$ : $\Gamma^p(\Delta)
(u_p) = \sum_{i+j=p} u_i \otimes u_j$.

Les « deux questions » de la page 323 : $\operatorname{Hom}_k(A, J)$ est-il un
idéal à puissances divisées de l'anneau $\operatorname{Hom}_k(A, B)$ (c'est
$\check{A} \otimes J$ si $A$ est libre de type fini) ? Les $u_n = u_1^{(n)}$
pour $B = \Gamma^*A$ satisfont-ils la relation, et sont-ils les seuls ? Et
dualement pour une algèbre augmentée et une cogèbre à coproduits divisés. Ces
pages, et les suivantes de même sujet, sont des calculs sans conclusion
énoncée.

\pagerange{325}{327}
\section*{Un feuillet étranger : la structure projective d'une courbe (pages 325 à 327)}

Trois feuillets à l'encre bleue, écrits en travers, changent de sujet. Pour une
courbe $X$ projective lisse connexe sur $\mathbb{C}$ de genre $g \geq 2$,
l'uniformisation donne un torseur canonique sous $GP(1)_{\mathbb{C}}$ muni d'une
connexion intégrable : en chaque $x$, les isomorphismes du revêtement universel
pointé avec le demi-plan de Poincaré forment un torseur sous $\operatorname{Aut}
(D_0) \simeq GP(1, \mathbb{R})$, d'où par extension un torseur sous
$GP(1,\mathbb{C})$, plat, associé à $\pi_1 \to GP(1, \mathbb{R}) \subset
GP(1,\mathbb{C})$\footnote{$GP(1)$ est sa notation pour le groupe projectif
$PGL_2$, quotient de $SL_2$ par $\mu_2$, comme le dit la suite exacte de la page
325. Un énoncé « une et une seule application analytique » de $\tilde{X}(x)$
dans $T_x(X)$, d'image un disque et de dérivée l'identité, y est noté.}. La
question encadrée : trouver une connexion intégrable algébrique canonique sur
$\mathbb{P}(\Omega^{-1} \oplus \Omega) = \overline{V(\Omega^2)}$. La page 327
calcule les classes d'Euler et l'obstruction à relever le torseur de
$GP(1,\mathbb{R})$ à $SL(2,\mathbb{R})$ dans $H^2(X, \mathbb{Z}/2)$, avec $E' =
\frac12 E$ et « $2g - 2 \sim g - 1$ »\footnote{Les diagrammes de la page 327
relient $H^4(B_{SL_2}, \mathbb{Z})$, $H^2(B_{GP(1)}, \mathbb{Z})$ et leurs
réductions modulo $2$ ; nous ne les restituons pas au-delà de ce qu'en dit ce
résumé. Aujourd'hui on parle de \emph{structure projective} (ou d'oper) sur la
courbe, et de l'inégalité de Milnor–Wood pour le nombre d'Euler.}. Ces calculs
sont épars et la lecture en est mince.

\pagerange{328}{334}
\section*{Gribouillis, suite : la structure à puissances divisées induite (pages 328 à 334)}

Retour au sujet des pages 319 à 324. Le foncteur $F_A$ commute au produit
tensoriel ; on définit $u^X_* : \Gamma^*F(X) \to F(\Gamma^*X)$ par $x^{(p)} =
u^X_p(\gamma^p(x))$, et l'on vérifie que c'est un homomorphisme d'algèbres et que
$x \mapsto x^{(p)}$ satisfait l'additivité, $(\lambda x)^{(p)} =
\lambda^px^{(p)}$, $x^{(p)}x^{(q)} = c_{p,q}x^{(p+q)}$, et — par le diagramme de
la page 331 — $(x^{(p)})^{(q)} = \frac{(pq)!}{(p!)^q\,q!}\,x^{(pq)}$. La page 332
réduit la compatibilité d'un morphisme aux puissances divisées à la vérification
sur des générateurs ; la page 333 écrit la liste des axiomes d'une algèbre $A$
opérant sur un idéal $J$ « sans $1$ » muni de $\gamma^p : J \to J$, et se demande
si $F(J) \subset F(A)$ hérite des puissances divisées. La moitié inférieure de
la page 329, qui définit $\operatorname{Sym}^*A \to A$ par $x_1\cdots x_n \mapsto
\frac{1}{n!}x_1\cdots x_n$, est barrée ; la page 330 ne porte qu'une ligne
biffée.

\pagerange{336}{340}
\section*{Gribouillis : rangs et cochaînes normalisées (pages 336 à 340)}

La page 336 compte les rangs, sur $\Delta^N$ : $\Phi_{(N)}$ est de rang $N+1$,
$\Psi_{(N)}$ de rang $N$, et
\[
\operatorname{rg}\bigl(\Gamma^p\Phi_{(N)} \otimes \Lambda^q\Psi_{(N)}\bigr) = \binom{N+p}{p}\binom{N}{q} = c_{p,q}\binom{N+p}{p+q} ,
\]
avec des relations de Pascal itérées $\tau^p\varphi_r = \sum_j \binom{p}{j}
\varphi_{r-j}$ sur les polynômes de rang, pour comparer aux $\Lambda^{p+q}\Phi_*$
— des cochaînes normalisées.

Les pages 337 à 340 dressent le tableau des quatre avatars des cochaînes
normalisées du simplexe : $\mathbb{C}^m_* = \Lambda^{m+1}\Phi_*$, avec
$\mathbb{C}^m_{[n]} = \Lambda^{m+1}k^{n+1} = k^{\operatorname{Inj}_\Delta(\Delta^m,
\Delta^n)}$, nul si $m > n$ et égal à $k$ si $m = n$ ; le complexe normalisé
$\mathbb{C}^m_\bullet$ est $k$ en degrés $m$ et $m+1$ ; et les duaux. Pour
$X_*$, $\operatorname{Hom}(X_*, \mathbb{C}^i_*) = \mathbb{C}^i(X_*)$ (cochaînes
non dégénérées), $\mathbb{C}_*(X_*) = k^{(X_*)}$, $\mathbb{C}^*(X_*) \simeq
\mathbb{C}_*(X_*)^\vee$. À $X_*$ il associe le foncteur $F_{X_*} : M_* \mapsto
\operatorname{Hom}(X_*, M_*)$, et calcule $\operatorname{Hom}(F_{X_*},
F_{\Delta^n_*}) \simeq k[X_{[n]}]^{\vee\vee}$, isomorphe à $k[X_{[n]}]$ si
$X_{[n]}$ est fini. Il note que $k^{\operatorname{Hom}(\Delta_n, \Delta_m)} \to
\operatorname{Hom}_k((k^{\Delta_m}, T_m), (k^{\Delta_n}, T_n))$ « n'est pas
bijectif ! ». La page 340, numérisée à l'envers, construit l'objet cosimplicial
$n \mapsto C^{[n]}_\bullet$ de $\mathrm{Ab}_{k\bullet}$ et son image par
$\operatorname{Hom}_\bullet(-, \tilde{\lambda}_\bullet)$\footnote{Sur la page
340 la seconde condition du dernier tableau répète « $j \neq i, i+1$ » pour
« $j = i, i+1$ ». La phrase sur l'objet cosimplicial est en grande partie
incertaine.}.

\pagerange{341}{345}
\section*{Retrouver $X$ à partir de ses cochaînes ; l'algèbre des opérateurs (pages 341 à 345)}

Pour un ensemble simplicial $X_*$, les cochaînes normalisées $C^\bullet(X_*)$ ont
une structure multiplicative ; il pose $F_*(X_*) = \bigl(n \mapsto
\operatorname{Hom}(C^\bullet(X_*), C^\bullet(\Delta^n_*))\bigr)$ et se demande si
$X_* \to F_*(X_*)$ est un isomorphisme, ou une équivalence d'homotopie (« est-ce
une équiv.\ d'homotopie ? », page 344) : on a $\operatorname{Hom}(\Delta^m_*,
\Delta^n_*) \xrightarrow{\sim} \operatorname{Hom}_{\mathrm{mult}}(F(\Delta^n_*),
F(\Delta^m_*))$, mais $F_*(\Delta^m) = P^m_*$ est un torseur sous
$\Psi_m^\vee \otimes \Phi_*$, et l'échange de $\varprojlim$ et
$\operatorname{Hom}$ est en question. L'exemple du cercle $X_* =
\operatorname{Coker}(\Delta^0_* \rightrightarrows \Delta^1_*)$ est commencé. La
page 341 introduit un foncteur $K$, noyau de $M_{[1]} \rightrightarrows M_{[0]}$,
et sa compatibilité au produit tensoriel. La page 342 contient un argument
d'indépendance linéaire de caractères (homomorphismes $u_i$ distincts, idéaux
premiers $\ker u_i$) dont la conclusion ne se lit pas.

La page 345, feuillet couvert de calculs serrés, dresse l'algèbre des
opérateurs agissant sur $\Gamma\Phi \otimes \Lambda\Psi$ : $d$, la multiplication
$\lambda$ par $\overline{T} = dT$, l'opérateur de Koszul $\delta$, les
multiplications $\mu_f$ par les $f \in k\{T\}$, et l'opérateur de degré
$\Theta(x) = \deg(x)\,x$. Relations :
\[
d\delta + \delta d = \Theta, \quad \delta\lambda + \lambda\delta = \mu_T, \quad d\mu_f - \mu_fd = \mu_{f'}\lambda, \quad \Theta\lambda - \lambda\Theta = \lambda, \quad \Theta\mu_f - \mu_f\Theta = \mu_{\Theta f} ,
\]
et $d^2 = \delta^2 = \lambda^2 = 0$, $d\lambda + \lambda d = 0$, $\lambda\mu_f =
\mu_f\lambda$, $\delta\mu_f = \mu_f\delta$, $\Theta$ commutant à $d$ et
$\delta$ ; tout élément de l'algèbre engendrée s'écrit $\sum \mu_f
\lambda^\alpha d^\beta \delta^\sigma \Theta^n$, $\alpha, \beta, \sigma \in
\{0,1\}$\footnote{La relation $\delta\lambda + \lambda\delta = \mu_T$ suppose
$\delta(\overline{T}) = T$ ; la page écrit « $\delta(\overline{T}) = 1$ », lecture
incertaine. C'est une algèbre du type des relations de Cartan ($d\iota +
\iota d = \mathcal{L}$) de la géométrie différentielle.}. Dans la colonne de
droite, un début de construction du scindage de la page 77 pour $q = 1$ et $p$
quelconque, par $(\varphi, \psi) \mapsto s(\varphi \otimes \psi)\,
\varphi^{(p-1)}$.

\pagerange{347}{351}
\section*{Algèbres différentielles sur un topos (pages 347 à 351)}

Trois pages au crayon très pâle, qui ne se lisent que par fragments : la lecture
est mince ici. Pour un topos $\mathcal{X}$ et une résolution $k \to C^*$, il
cherche à munir $C^*$ d'« accouplements » partiels $Z^i \times C^j \to C^{i+j}$
et $C^i \times Z^j \to C^{i+j}$, coïncidant sur $Z^i \times Z^j$ à valeurs dans
les cocycles, associatifs, unitaires, anticommutatifs, avec $x \circ x = 0$
pour $x$ cocycle de degré impair, et compatibles à $d$ ; sous ces conditions,
$C^*(U) = \Gamma(U, C^*_U)$ est une $k(U)$-algèbre et $H^*(C^*(U)) \to H^*(U,
k_U)$ est un homomorphisme d'algèbres. L'exemple est $C^i = \Lambda^{i+1}\Phi$,
de cocycles $Z^i \simeq \Lambda^i\Psi$, et la page 351 calcule
$\operatorname{Hom}_k(C^i \otimes Z^j, C^{i+j}) \simeq \Lambda^i\Psi \otimes
\Lambda^j\Psi \simeq \Lambda^{i+j}\Psi \oplus (\cdots)$. Les pages 348, 350 et
352 sont un tapuscrit d'une autre main.

\pagerange{354}{354}
\section*{Le résumé des foncteurs universels (page 354)}

Feuillet qui condense les pages 310 à 317 : pour un type $\mathcal{L}$ de
limites, $\underline{\operatorname{Hom}}^!_{\mathcal{L}}(\mathcal{S}s^\circ,
\mathcal{B}) \simeq \underline{\operatorname{Hom}}^!_{\mathcal{L}}(\mathcal{B}_0,
\mathcal{B})$ avec $\mathcal{B}_0 \simeq (\mathrm{Ab}_*)^\circ$, et, encadré,
$F(X_*) = \Omega_*(\mathbb{C}_*(X)) = \Omega_\bullet(\mathbb{C}_\bullet(X))$ ;
la moitié tournée de la page refait l'argument de représentabilité des pages 312
et 313.

\pagerange{356}{360}
\section*{Notes : Eilenberg--Zilber et les opérations de Steenrod (pages 356 à 360)}

Les pages 356 à 375 sont d'une autre facture que le reste : des notes, avec
références (Mac~Lane, \emph{Homology} ; Steenrod–Epstein, \emph{Cohomology
operations} ; la thèse d'Illusie ; Dold ; Cartan ; Serre ; Lazard ; Milne,
lecture incertaine), qui ressemblent à des notes d'exposés ou de lecture. Les
pages ne disent pas de qui elles sont, et nous ne l'attribuons
pas\footnote{Rien sur les feuillets ne permet de trancher entre des notes
prises par Grothendieck à des exposés et des notes de sa propre lecture.}.

\textbf{Théorème d'Eilenberg–Zilber–Cartier.} \emph{Pour un objet abélien
bisimplicial $X_{**}$, il existe des morphismes naturels $\varphi : \int
N X_{**} \to N\Delta X_{**}$ et $\psi$ en sens inverse, entre le complexe total
du bicomplexe normalisé et le normalisé de la diagonale, qui coïncident avec
les flèches évidentes en degré $0$ et sont inverses l'un de l'autre à homotopie
naturelle près ; les homotopies sont elles-mêmes uniques à homotopie près.} Les
choix traditionnels sont l'application d'Alexander–Whitney et celle des
\emph{shuffles}, toutes deux associatives, la seconde (anti)commutative, la
première commutative seulement à homotopie près ; on peut l'exprimer par des
objets trisimpliciaux (thèse d'Illusie). E–Z–C entraîne E–Z, et le
cup-produit $h = \sum \pm f(x_1,\ldots,x_i)g(x_{i+1},\ldots)$ est associatif et
commutatif à homotopie près, d'où l'anticommutativité de $H^*(X, R)$. Les
produits $K(\mathbb{Z},i) \wedge K(\mathbb{Z}, j) \to K(\mathbb{Z}, i+j)$ se
construisent par $N(K \otimes K) \to NK \otimes NK = \mathbb{Z}[i] \otimes
\mathbb{Z}[j]$.

\emph{Corollaire.} Avec $W_n = \mathbb{Z}E\Sigma_n$, résolution libre de
$\mathbb{Z}$ par des $\Sigma_n$-modules, on a des morphismes $\Sigma_n$-équivariants
$W_n \otimes \mathbb{Z}(X^n) \to (\mathbb{Z}X)^{\otimes n}$ et inverses, qui
expriment la commutativité à homotopie près, avec toutes ses cohérences — comme
pour un H-espace $E\Sigma_n \times X^n \to X$ relevant la multiplication itérée,
$E\Sigma_2 = \varinjlim S^n$. \textbf{Théorème} (Dold, « Über Steenrodsche
Operationen ») : $\operatorname{Hom}^\bullet(\mathbb{Z}(X^n), (\mathbb{Z}X)^{\otimes
n})$ est, pour $X$ variable, une résolution libre de $\mathbb{Z}$ dans les
$\Sigma_n$-modules.

\emph{Construction des opérations de Steenrod} (cf.\ Steenrod–Epstein, p.~99).
Pour $A$ de caractéristique $p$ et $f \in H^n(X, A)$, le composé
\[
W_p \otimes \mathbb{Z}(X^p) \xrightarrow{\ \lambda\ } (\mathbb{Z}X)^{\otimes p} \xrightarrow{\ f^{\otimes p}\ } A[n]^{\otimes p} \xrightarrow{\ \mu\ } A[pn]
\]
donne $Pf \in H^{pn}_{\Sigma_p}(E\Sigma_p \times X^p, A)$, et par la diagonale
$\Delta Pf \in H^{pn}(K(\Sigma_p,1) \times X, A) = H^*(K(\Sigma_p,1))
\otimes_{\mathbb{F}_p} H^*(X, A)$ ; les composantes $D_ix \in H^{pn-i}(X,A)$
selon une base $e_\alpha$ de $H_*(K(\Sigma_p, 1))$ donnent, réindexées, les
$P^i$, et « ce réindexage est ce qui donne la formule de Cartan ». On travaille
en fait avec $\mathbb{Z}/p \subset \Sigma_p$.

\pagerange{361}{367}
\section*{Notes : décalage, espaces d'Eilenberg--MacLane, algèbre de Steenrod (pages 361 à 367)}

\emph{Formule de décalage (Quillen–Bousfield, cf.\ thèse d'Illusie).} Pour les
foncteurs dérivés au sens de Dold–Puppe,
\[
\mathrm{L}\operatorname{Sym}^i(X[1]) \simeq \mathrm{L}\Lambda^i(X)[i] , \qquad \mathrm{L}\Lambda^i(X[1]) \simeq \mathrm{L}\Gamma^i(X)[i] , \qquad \mathrm{L}\operatorname{Sym}^i(X[2]) \simeq \mathrm{L}\Gamma^i(X)[2i] ,
\]
à partir des suites de Koszul $0 \to \operatorname{Sym}^iK(M,n) \to \cdots \to
\Lambda^iE(M,n+1) \to \Lambda^iK(M,n+1) \to 0$ associées à $K(M,n) \to E(M,n+1)
\to K(M,n+1)$ ; encadré : $\mathrm{L}_r\operatorname{Sym}^iK(M,n+1) \simeq
\mathrm{L}_{r-i}\Lambda^iK(M,n)$\footnote{La page note les décalages $[-i]$ pour
ce que nous écrivons $[i]$ (déplacement vers le haut en homologie). Les deux
premières lignes de la page, en indices $L_i\operatorname{Sym}^i(M,n) \simeq
L_{i+n}\Lambda^i(M,n+1)$, sont traversées de traits obliques et ne s'accordent
pas avec l'énoncé encadré, que nous suivons.}.

\emph{Dold–Thom et Dold–Puppe.} $\pi_i\operatorname{SP}^\infty X \simeq H_i(X)$ ;
le théorème de Dold $\mathbb{Z}\operatorname{SP}^\infty X \simeq
\operatorname{Sym}\mathbb{Z}X$ donne $H_*(K(G,n)) = \mathrm{L}_*\operatorname{Sym}(G,
n)$, d'où une graduation « par le poids » sur l'homologie des espaces
d'Eilenberg–MacLane, et la stabilité $H_{n+r}K(M,n) \simeq H_{n+r+1}K(M,n+1)$
pour $r < n$.

\emph{Opérations stables.} Pour $r < n$, $H^{n+r}(K(\mathbb{F}_p, n),
\mathbb{F}_p)$ est stable ; les opérations stables forment l'algèbre de
Steenrod, $\mathbb{F}_p[P^0, P^1, \ldots, \beta]/(\text{Adem}, P^0 = 1)$, $\deg P^i
= 2i(p-1)$, $\deg\beta = 1$, pour $p \neq 2$, et $\mathbb{F}_2[\mathrm{Sq}^i]/
(\text{Adem}, \mathrm{Sq}^0 = 1)$, $\deg\mathrm{Sq}^i = i$, avec $\mathrm{Sq}^1 =
\beta$, pour $p = 2$ — par exemple $\mathrm{Sq}^1\mathrm{Sq}^2 =
\mathrm{Sq}^3$\footnote{La page écrit « $\mathrm{Sq}^1\mathrm{Sq}^2 =
\mathrm{Sq}^3 + \mathrm{Sq}^2\mathrm{Sq}^1$ (?) », avec le point d'interrogation ;
la relation d'Adem donne $\mathrm{Sq}^1\mathrm{Sq}^2 = \mathrm{Sq}^3$ (et
$\mathrm{Sq}^1\mathrm{Sq}^1 = 0$).}. Formule de Cartan $P^i(x \times y) = \sum_j
P^j(x)P^{i-j}(y)$ ; monômes $P^I = \beta^{\varepsilon_0}P^{\eta_1}
\beta^{\varepsilon_1}\cdots$, admissibles si $\eta_i \geq p\eta_{i+1} +
\varepsilon_i$, qui forment une base ; excès $e(I)$\footnote{La page donne pour
l'excès la forme « $\eta_1 - \sum_{i>1}\eta_i$ », qui est celle du cas $p = 2$ ;
pour $p$ impair l'excès compte aussi les Bockstein.}. \emph{Théorème} (Serre
pour $p = 2$, Cartan) : $H^*(K(\mathbb{F}_p, n), \mathbb{F}_p)$ est l'algèbre
commutative graduée libre engendrée par les $P^I\iota_n$, $I$ admissible
d'excès $< n$ (pour $p$ impair, s'y ajoutent ceux d'excès $n$ « qui se terminent
par $\beta$ ») ; l'homologie est l'algèbre à puissances divisées sur les
éléments duaux. Démonstrations : par la suite spectrale du classifiant
(séminaire Cartan, bar construction, Eilenberg–Moore, Rothenberg–Steenrod), ou
par la méthode de Serre (1954).

\emph{L'algèbre duale de Milnor.} La formule de Cartan donne la comultiplication
$\Delta P^i = \sum P^j \otimes P^k$ ; le dual est $A_* = \operatorname{Sym}(\xi_i)
\otimes \Lambda(\tau_j)$, $\deg\xi_i = 2p^i - 2$, $\deg\tau_j = 2p^j - 1$, avec
$\Delta\xi_k = \sum_i \xi_{k-i}^{p^i} \otimes \xi_i$ et $\Delta\tau_k = \tau_k
\otimes 1 + \sum_i \xi_{k-i}^{p^i} \otimes \tau_i$. Le cup-produit $K(\mathbb{F}_p,
n) \wedge K(\mathbb{F}_p, m) \to K(\mathbb{F}_p, m+n)$ induit la comultiplication
de $A$. $\mathrm{L}\operatorname{Sym}^r = 0$ pour $r \neq p^s$, et
$\mathrm{L}^{\mathrm{st}}_*\operatorname{Sym}\,\mathbb{F}_p = A_*$. La page 367
rappelle la bar construction $\overline{B}(A)$ d'une algèbre différentielle
graduée augmentée, $\overline{B}(\mathbb{Z}K(G,n)) \sim \mathbb{Z}BK(G,n)$, et sa
structure d'algèbre via Eilenberg–Zilber quand $A$ est commutative.

\pagerange{369}{372}
\section*{Notes : $K_2$ de Milnor ; cohomologie des Eilenberg--MacLane (pages 369 à 372)}

\emph{$K_2$ de Milnor.} $E(R) = [E(R), E(R)] \subset GL(R)$, $K_1(R) = GL(R)/E(R)$,
et $E(R) = SL(R)$ si $R$ est local ; la suite $0 \to K_2(R) \to \mathrm{St}(R) \to
E(R) \to 0$, et $K_2(R) = H_2(E(R), \mathbb{Z}) = \pi_2(BGL(R)^+)$\footnote{La page
écrit « $\pi_1(E(R)) = K_2(R)$ », l'indice étant incertain ; $E(R)$ est un
groupe discret, et c'est $H_2(E(R))$, ou $\pi_2$ de $BGL(R)^+$ qui figure à la
ligne suivante, qui donne $K_2$.}. Pour un schéma $X$, la suite spectrale
$H^i_{\mathrm{Zar}}(X, \underline{K}_j) \Rightarrow K_{j-i}(X)$, la résolution de
Gersten $0 \to \underline{K}_i \to i_*\underline{K}_i \to \bigoplus_{\mathrm{codim}\,1}
i_{x*}K_{i-1} \to \cdots \to \bigoplus_{\text{points fermés}} i_{x*}\mathbb{Z}
\to 0$, et $H^i(X, \underline{K}_i) = $ cycles de codimension $i$ modulo
l'équivalence rationnelle\footnote{La page écrit « $i$-cycles » ; c'est la formule
de Bloch–Quillen, en codimension $i$, et la suite de Gersten s'écrit pour $X$
régulier.}.

\emph{Cohomologie des Eilenberg–MacLane.} Les quatre cas $H_*(K(\mathbb{F}_p, n),
\mathbb{F}_p)$, $H_*(K(\mathbb{Z},n), \mathbb{F}_p)$, $H_*(K(\mathbb{F}_p,n),
\mathbb{Z})$, $H_*(K(\mathbb{Z},n),\mathbb{Z})$ ; pour $G$ quelconque, la partie
stable de l'homologie entière est de torsion première, par exemple
$H_{n+2}(K(G,n)) = G/2G$, $H_{n+3} = {}_2G$. Les deux Bockstein $\beta_1$, $\beta_2$
sont des dérivations de $A_* = \operatorname{Sym}(\xi_i) \otimes \Lambda(\tau_j)$,
$\tau_0 \mapsto 1$ pour l'un, $\tau_j \mapsto \xi_j$ pour l'autre ; d'où
$H_*(K(\mathbb{F}_p, n), \mathbb{Z})$ et $H_*(K(\mathbb{Z},n), \mathbb{Z})_{(p)}
= \operatorname{Im}\beta_2 = \operatorname{Ker}\beta_2$ dans le domaine stable,
et la conjugaison de Steenrod échange les deux Bockstein.

\emph{Théorème (Cartan)} : sous des hypothèses sur les groupes $G$, $H$,
l'homologie stable est symétrique : $H_{m+i}(K(H,m); G) \simeq H_{n+i}(K(G,n);
H)$ dans le domaine stable\footnote{La page écrit « $H_*(K(G,n),H) \simeq
H_*(K(H,n),G)$ additif », sans les décalages ; la seconde démonstration
($\pi_{i+m}(K(H,m) \wedge K(G,n)) \simeq \pi_{i+n}(K(G,n) \wedge K(H,m))$ par
permutation des facteurs) fixe les degrés que nous écrivons.}. Démonstration
2 : $X \mapsto \pi_{n+r}(K(\mathbb{Z},n) \wedge X)$ est une théorie
d'homologie partielle pour $r$ petit, égale à $\widetilde{H}_r(X)$, via
$K(\mathbb{Z},n) \wedge X \to \mathbb{Z}^+(S^nX)$, quasi-isomorphisme jusqu'en
degré $2n-1$, et Eilenberg–Zilber pointé. La page 372 en déduit des
résolutions simpliciales canoniques de $K(G,n) = G \otimes \Lambda^n\Psi$ et la
connectivité $r = \min(n + 2p - 1, 2n + p - 1)$ de $K(\mathbb{Z},n) \wedge K(G,p)
\to \mathbb{Z}K(G,n+p)$.

\pagerange{373}{375}
\section*{Notes : les opérations sur $\mathbf{G}_a$ et $\operatorname{Ext}(\mathbf{G}_a, \mathbf{G}_a)$ (pages 373 à 375)}

Pour un anneau $R$ de caractéristique $p$, $H^q(X, R) \simeq H^q(X, \mathbb{F}_p)
\otimes R$, et l'algèbre de Steenrod agit avec $P^0_R = P^0 \otimes F$ ($F$, le
Frobenius, lecture probable). Sur le site des schémas sur $\mathbb{F}_p$, avec le
faisceau structural $\mathbf{G}_a$ et $K(\mathbf{G}_a, n) = \mathbf{G}_a \otimes
\Lambda^n\Psi$, l'hypercohomologie $\mathbb{H}^*(K(F,p), G)$ classe les
opérations de type $(F,p;G,q)$ (Illusie ; Epstein, « Operations in homological
algebra », 1966). \emph{Théorème} : soit $\mathcal{A}$ l'algèbre engendrée par les
$P^i$, $i \geq 0$, avec les relations d'Adem \emph{mais pas} $P^0 = 1$ ; alors
$\mathbb{H}^*(K(\mathbf{G}_a, n), \mathbf{G}_a)$ est l'algèbre commutative graduée
libre engendrée par les monômes admissibles de $\mathcal{A}$ d'excès $< n$\footnote{La
page écrit $K(\mathcal{O}, m)$ et « excès $< n$ » ; nous uniformisons. Le calcul
de $H^*(K(\mathbf{G}_a, 1), \mathbf{G}_a)$ est renvoyé à Lazard (« Lois de groupes
et analyseurs ») ; on grimpe ensuite par le théorème de Borel et la suite de
Leray de $K(\mathbf{G}_a, n-1) \to * \to K(\mathbf{G}_a, n)$. La classe $\beta \in
H^2(K(\mathbf{G}_a,1), \mathbf{G}_a)$ est celle de l'extension de Witt $0 \to
\mathbf{G}_a \to W_2 \to \mathbf{G}_a \to 0$.}.

\emph{Proposition} : le $k[F]$-module $\operatorname{Ext}^i(\mathbf{G}_a,
\mathbf{G}_a)$ est de $F$-torsion pour $i > 1$. \emph{Corollaire} (attribué,
lecture incertaine, à Milne) : sur le site parfait, avec $\mathbf{G}_a^{\mathrm{parf}}
= \varinjlim_F \mathbf{G}_a$, $\operatorname{Ext}^i(\mathbf{G}_a^{\mathrm{parf}},
\mathbf{G}_a^{\mathrm{parf}}) = 0$ pour $i > 1$. La démonstration utilise la suite
spectrale du coefficient universel $\operatorname{Ext}^i(\underline{H}_j(X), F)
\Rightarrow \mathbb{H}^{i+j}(X, F)$ pour $X = K(\mathbf{G}_a, n)$, $n$ grand, dans
les faisceaux de $\mathbb{F}_p$-modules et « modulo $F$ », avec $\operatorname{Hom}
(\mathbf{G}_a, \mathbf{G}_a) = k[F]$ et $\operatorname{Ext}^1(\mathbf{G}_a, \mathbf{G}_a)
= 0$\footnote{Cette dernière annulation est à lire dans la catégorie des faisceaux
de $\mathbb{F}_p$-modules, où la page dit travailler : l'extension de Witt $W_2$,
qui n'est pas tuée par $p$, n'y vit pas, et ne la contredit donc pas.}, deux
lemmes (nullité d'un homomorphisme de bord ; générateurs $P^I\iota_n$), et la
résolution par les $\mathbb{Z}^jK(\mathbf{G}_a, n)$. La page 375 inscrit le
résultat : $\operatorname{Ext}^{2i+1}_{\mathbb{F}_p}(\mathbf{G}_a, \mathbf{G}_a) = 0$
et $\operatorname{Ext}^{2i}(\mathbf{G}_a, \mathbf{G}_a) = k[F]/(F^d)$, $d =
v_p(i)$, pour $i > 0$\footnote{Nous le restituons tel que la page l'écrit, sans
l'avoir vérifié. Ces calculs sont aujourd'hui associés au nom de L.~Breen, qui
publiait « Extensions du groupe additif » dans les \emph{Publications
mathématiques de l'IHÉS} en 1978 ; les pages ne le nomment pas, et nous ne
prétendons pas savoir de quel exposé elles procèdent. L'argument s'interrompt
avec le dossier.}.

\end{document}
